\documentclass[11pt]{article}

\usepackage[margin=1in]{geometry}
\usepackage[T1]{fontenc}
\usepackage[utf8]{inputenc}
\usepackage{lmodern}
\usepackage{microtype}
\usepackage{authblk}
\usepackage{amsmath,amssymb,amsthm,mathtools}
\usepackage{bm}
\usepackage{xcolor}
\usepackage{enumitem}
\usepackage{aliascnt}
\usepackage{array,booktabs,multirow,tabularx}
\newcounter{algorithm}
\newcommand{\SetKwInput}[2]{}
\newcommand{\SetKwComment}[3]{}
\newcommand{\SetKw}[2]{}
\newcommand{\KwInput}[1]{\noindent\textbf{Input: }#1\par}
\newcommand{\KwOutput}[1]{\noindent\textbf{Output: }#1\par}
\newcommand{\tcp}[1]{}
\newcommand{\Return}{\textbf{return} }
\newcommand{\While}[2]{%
  \noindent\textbf{while }#1\textbf{ do}\par
  \begingroup\leftskip=1.5em #2\par\endgroup
  \noindent\textbf{end while}\par}
\newcommand{\If}[2]{%
  \noindent\textbf{if }#1\textbf{ then }#2\par}
\newenvironment{algorithm}[1][]{%
  \par\medskip\noindent\begin{minipage}{\linewidth}\small
  \hrule\smallskip
  \renewcommand{\caption}[1]{%
    \refstepcounter{algorithm}%
    \noindent\textbf{Algorithm~\thealgorithm: ##1}\par\smallskip}
  \renewcommand{\;}{\par}%
}{%
  \smallskip\hrule\end{minipage}\par\medskip}
\usepackage{hyperref}
\usepackage[capitalize,nameinlink]{cleveref}
\crefname{algorithm}{algorithm}{algorithms}
\Crefname{algorithm}{Algorithm}{Algorithms}

\allowdisplaybreaks
\hypersetup{
  colorlinks=true,
  linkcolor=blue!55!black,
  citecolor=green!40!black,
  urlcolor=blue!55!black
}

\newtheorem{theorem}{Theorem}[section]
\newaliascnt{lemma}{theorem}
\newtheorem{lemma}[lemma]{Lemma}
\aliascntresetthe{lemma}
\newaliascnt{corollary}{theorem}
\newtheorem{corollary}[corollary]{Corollary}
\aliascntresetthe{corollary}
\newaliascnt{proposition}{theorem}
\newtheorem{proposition}[proposition]{Proposition}
\aliascntresetthe{proposition}
\newaliascnt{definition}{theorem}
\newtheorem{definition}[definition]{Definition}
\aliascntresetthe{definition}
\newaliascnt{remark}{theorem}
\newtheorem{remark}[remark]{Remark}
\aliascntresetthe{remark}
\newaliascnt{fact}{theorem}
\newtheorem{fact}[fact]{Fact}
\aliascntresetthe{fact}

\newcommand{\bits}{\{0,1\}}
\newcommand{\cA}{\mathcal{A}}
\newcommand{\cB}{\mathcal{B}}
\newcommand{\cD}{\mathcal{D}}
\newcommand{\cE}{\mathcal{E}}
\newcommand{\cH}{\mathcal{H}}
\newcommand{\cK}{\mathcal{K}}
\newcommand{\cM}{\mathcal{M}}
\newcommand{\cR}{\mathcal{R}}
\newcommand{\cS}{\mathcal{S}}
\newcommand{\cU}{\mathcal{U}}
\newcommand{\cX}{\mathcal{X}}
\newcommand{\cY}{\mathcal{Y}}
\newcommand{\Dom}{\operatorname{Dom}}
\newcommand{\Unif}{\operatorname{Unif}}
\newcommand{\poly}{\operatorname{poly}}
\newcommand{\negl}{\operatorname{negl}}
\newcommand{\Tr}{\operatorname{Tr}}
\newcommand{\supp}{\operatorname{supp}}
\newcommand{\Supp}{\operatorname{Supp}}
\newcommand{\Real}{\mathsf{Real}}
\newcommand{\Learn}{\mathsf{Learn}}
\newcommand{\Heavy}{\mathsf{Heavy}}
\newcommand{\View}{\mathsf{View}}
\newcommand{\IntKR}{\mathsf{IntKR}}
\newcommand{\Adv}{\mathsf{Adv}}
\newcommand{\Gen}{\mathsf{Gen}}
\newcommand{\Enc}{\mathsf{Enc}}
\newcommand{\Dec}{\mathsf{Dec}}
\newcommand{\IND}{\mathsf{IND}}
\newcommand{\pk}{\mathsf{pk}}
\newcommand{\sk}{\mathsf{sk}}
\newcommand{\ct}{\mathsf{ct}}
\newcommand{\ket}[1]{\lvert #1\rangle}
\newcommand{\bra}[1]{\langle #1\rvert}
\newcommand{\proj}[1]{\ket{#1}\!\bra{#1}}
\newcommand{\given}{\,\middle|\,}
\newcommand{\ceil}[1]{\left\lceil #1\right\rceil}
\newcommand{\changed}[1]{\textcolor{purple!75!black}{#1}}
\newcommand{\mrchanged}[1]{\changed{#1}}

\newenvironment{hybrid}[1]{%
  \par\medskip
  \noindent\textbf{#1}
  \begin{enumerate}[leftmargin=2.5em,label=\arabic*.,itemsep=0.25em,topsep=0.35em]
}{%
  \end{enumerate}
  \medskip
}

\SetKwInput{KwInput}{Input}
\SetKwInput{KwOutput}{Output}
\SetKwComment{tcp}{\(\triangleright\) }{}
\SetKw{Return}{return}

\title{Towards the Impossibility of Imperfectly Complete Key Agreement in the QROM}
\renewcommand{\Authfont}{\normalsize}
\renewcommand{\Affilfont}{\small}
\renewcommand{\Authands}{, }
\setlength{\affilsep}{0.45em}
\author[1]{Fuyuki Kitagawa\thanks{\texttt{fuyuki.kitagawa@ntt.com}}}
\author[1]{Ryo Nishimaki\thanks{\texttt{ryo.nishimaki@ntt.com}}}
\author[2]{Ági Villányi\thanks{\texttt{agivilla@mit.edu}}}
\author[1]{Takashi Yamakawa\thanks{\texttt{takashi.yamakawa@ntt.com}}}
\affil[1]{NTT, Tokyo, Japan}
\affil[2]{Massachusetts Institute of Technology, Cambridge, MA, USA}
\date{}


\begin{document}
\maketitle

\begin{abstract}
We make progress towards the impossibility of imperfectly complete
quantum-computation, classical-communication (QCCC) key agreement by
constructing the first
unconditional attacks on quantum key agreement in the following restricted settings. 

In the two-message setting, we assume that Alice makes only classical queries to the oracle in the first round and that her message to Bob is classical, but otherwise both parties may perform arbitrary quantum computation, make quantum queries, and send a quantum state in the second round. Our attack and analysis are based on the heavy-query learning techniques from Austrin et al. (CRYPTO 2022) and the reprogramming techniques of Katz and Sela (arXiv 2401.14319). 

In the round-independent setting, we show that the attack of Barak and Mahmoody (CRYPTO 2009; J. Cryptology 2017) can be extended to multiple rounds when Alice and Bob share classical communication and make only classical queries in all but the final round. In both settings, the attacker is computationally unbounded and makes \(\poly(\lambda)\) queries to recover the key whenever each
honest query bound is at most \(\poly(\lambda)\) and the valid agreement probability is
inverse-polynomial. As a consequence, we rule out imperfectly correct quantum
public-key encryption for classical messages whose length is bounded by a
polynomial in \(\lambda\)
in the QROM when key generation has classical oracle access, even if
encryption, decryption, and the ciphertext are quantum. In particular, the
one-bit case applies to the imperfectly correct PKE obtained from two-round
OSP by Bartusek and Khurana (CRYPTO 2025) whenever the classical
OSP sender makes only classical random-oracle queries. 
\end{abstract}

\tableofcontents

\section{Introduction}

Merkle studied key agreement as a way for two parties with no shared secret
to communicate securely over an insecure channel.  His puzzle protocol allows
the two parties to establish a key while requiring a passive eavesdropper to
make quadratically more queries to a random oracle~\cite{Merkle78}. Subsequent work showed
that, for classical protocols, this quadratic gap is optimal under both
perfect and imperfect completeness.  Impagliazzo and Rudich ruled out a
super-polynomial gap~\cite{ImpagliazzoRudich89}, and Barak and Mahmoody proved
a tight quadratic upper bound, including when the honest parties need not
agree with certainty~\cite{BarakMahmoodyFull}.

It is natural to ask whether these results can be extended to the quantum setting, in which parties are allowed to make queries to the oracle in superposition. This is known as the quantum random-oracle model (QROM)~\cite{BonehEtAl11} and quantum variants of Merkle's approach,
and key agreement in the QROM more broadly, have been studied in
several works~\cite{BrassardSalvail08,BrassardEtAl19}. The model we consider in this work was introduced by Austrin
et al.~\cite{AustrinEtAl22} as
quantum-computation, classical-communication (QCCC). In a QCCC
protocol, the parties may perform arbitrary quantum computation, retain
quantum states, and query the random oracle in superposition, but every
message they exchange is classical.

In 2018, Hosoyamada and Yamakawa asked whether QCCC key agreement in the
QROM can achieve a super-polynomial gap between
the query complexities of the honest parties and the
eavesdropper~\cite[Section~6]{HosoyamadaYamakawa18}. Assuming perfect completeness, Austrin et al. gave an unconditional quadratic-query classical attack when one party
is classical~\cite{AustrinEtAl22}. For general QCCC protocols, their Polynomial Compatibility
Conjecture yields a classical attack whose query complexity is polynomial in
the honest parties' query bounds; unconditionally, they obtained an attack using exponentially many queries. 

Li, Li, Li, and Liu subsequently obtained unconditional attacks in a sequence
of increasingly general, perfectly complete settings. They first treated
QPKE with classical-query key generation, while also proving a separate
imperfect-correctness result for oracle-free decryption~\cite{LiEtAl24}.
They then allowed logarithmically many quantum key-generation queries
unconditionally and formulated a conjecture for the general
case~\cite{LiEtAl25}. Their later work removed both restrictions for
perfectly complete QPKE with classical keys and either classical or quantum
ciphertexts~\cite{LiEtAl26PKE}. Most recently, they gave an unconditional,
round-independent attack on every perfectly complete QCCC key-agreement
protocol. The honest parties may query the oracle in superposition in every
round, while the eavesdropper makes only classical
queries~\cite{LiEtAl26ManyRound}.

What remains open is imperfectly complete QCCC key agreement when the honest
parties may make arbitrary quantum queries. Even a classical-query attack in
this setting would face a serious obstacle identified by Austrin et al. Their
Simulation Conjecture asserts, informally, that a classical algorithm whose
query complexity is polynomial in the oracle's input length can approximate
the acceptance probability of any quantum algorithm with the same type of
query bound, on most random oracles. It follows from
the Aaronson-Ambainis conjecture on influential variables of bounded
low-degree polynomials~\cite{AaronsonAmbainis14}. Austrin et al. showed that,
if the Simulation Conjecture is false, then there exists an imperfectly complete QCCC protocol secure against
every classical eavesdropper whose query complexity is polynomial in the
security parameter. A fully general classical
attack would therefore require substantial progress on this simulation
problem. The barrier leaves room, however, for quantum-query attacks and for
attacks that restrict the honest parties' quantum queries. We pursue both
possibilities in this work.

Before this work, no unconditional impossibility result was known for
imperfectly correct QCCC key agreement in the QROM with unrestricted quantum
oracle access, even for two-message protocols.  The two-message setting also
captures quantum public-key encryption (QPKE) through a standard
key-agreement protocol: Alice sends her public key, Bob responds with an
encryption of a uniformly random bit, and Alice recovers the shared bit by
decrypting. If the ciphertext is quantum, then the second message in this protocol is
quantum as well.  Bouaziz-Ermann et al. extended the conjectural approach to
this more general setting~\cite{BouazizEtAl24}.  Our two-message result gives
an unconditional attack when Alice's first-stage oracle queries and first
message are classical, and therefore addresses an open case not covered by
the perfectly complete impossibility results.

\subsection{Our Results}

We prove unconditional attacks that do not assume perfect completeness, where the attacker is a query-bounded, computationally unbounded party. 

For a key-agreement protocol \(\Pi\), let
\[
  \alpha_\Pi(\lambda)
  :=
  \Pr[K_{\cA}=K_{\cB}\neq\bot]
\]
denote its valid agreement probability. Our bounds are meaningful whenever
\(\alpha_\Pi\) is nonnegligible.  Thus they apply in a substantially broader
regime than negligible completeness error: the honest parties need only agree
on a valid key with inverse-polynomial probability, rather than with
probability \(1-\negl(\lambda)\). 

We begin with two-message protocols. Before sending the first message, Alice
queries the oracle only classically, although her remaining computation and
retained state may be quantum. Bob may return a quantum message, and both Bob
and Alice's final algorithm may query the oracle in superposition. Let
\(q_{\cA,1}\), \(q_{\cB}\), and \(q_{\cA,2}\) bound the number of queries
in Alice's first stage, Bob's stage, and Alice's final stage, respectively.

\begin{theorem}[Two-message attack, informal]
\label{thm:informal-two-message}
Let \(\Pi\) be a protocol of the form above with valid agreement probability
\(\alpha>0\).  There is a computationally unbounded quantum attacker that
recovers Bob's key with probability greater than \(\alpha/2\) using
\[
  O\!\left(
    \frac{
      q_{\cA,1}\max\{1,2q_{\cB}+q_{\cA,2}\}^{4}}
      {\alpha^{4}}
    +q_{\cA,2}
  \right)
\]
random-oracle queries.  Its queries before Bob's response are classical;
its queries after intercepting the response may be quantum.
\end{theorem}

Consequently, no family of such protocols with honest query bounds at most
\(\poly(\lambda)\) and inverse-polynomial valid agreement is secure against
computationally unbounded attackers making \(\poly(\lambda)\) queries. If Alice is
entirely classical and Bob's response is classical, the attack is a passive
classical eavesdropping attack.

In the multi-round setting, we allow for a classical-communication prefix of arbitrary length. The parties may
run quantum algorithms and store quantum states, but all oracle queries made before the final round must be classical. Bob next sends one possibly quantum message, and
Alice performs her final computation. These last two stages may query the
oracle in superposition. Let
\(q_{\cA,\mathsf{pre}}\) and \(q_{\cB,\mathsf{pre}}\) bound the parties'
total prefix queries, and let \(q_{\cB,\mathsf q}\) and
\(q_{\cA,\mathsf{fin}}\) bound their final-stage queries.

\begin{theorem}[Multi-round attack, informal]
\label{thm:informal-multi-round}
Let \(\Pi^{\mathsf{mr}}\) be a protocol of the form above with valid
agreement probability \(\alpha>0\).  There is a computationally unbounded
online quantum attacker that recovers Bob's key with probability greater
than \(\alpha/2\) using
\[
  O\!\left(
    \frac{
      \max\{1,q_{\cA,\mathsf{pre}}\}
      \max\{1,q_{\cB,\mathsf{pre}}\}
      \max\{1,2q_{\cB,\mathsf q}+q_{\cA,\mathsf{fin}}\}^{6}}
      {\alpha^{6}}
    +q_{\cA,\mathsf{fin}}
  \right)
\]
random-oracle queries.  Its queries during the classical prefix are
classical, and the bound has no additional dependence on the number of
prefix rounds.
\end{theorem}

As an application, we consider quantum public-key encryption for classical
messages whose length is bounded by a polynomial in \(\lambda\). The public
and secret keys are classical, and key generation has classical oracle
access. Encryption, decryption, and the ciphertext itself may be quantum.

\begin{theorem}[QPKE attack, informal]
\label{thm:informal-qpke}
No QPKE scheme for classical messages whose length is bounded by a polynomial
in \(\lambda\), with
classical-query key generation, negligible correctness error, and honest
query bounds at most \(\poly(\lambda)\), is query-bounded IND-CPA secure. In
fact, for every polynomial \(p\), there is a computationally unbounded quantum
attacker making \(\poly(\lambda)\) queries whose IND-CPA advantage is at least
\[
  \frac12-\frac{1}{p(\lambda)}-\negl(\lambda).
\]
\end{theorem}

\begin{table}[t]
  \centering
  \small
  \renewcommand{\arraystretch}{1.15}
  \begin{tabularx}{\textwidth}{@{}l*{3}{>{\centering\arraybackslash}p{1.45cm}}>{\centering\arraybackslash}p{1.75cm}>{\raggedright\arraybackslash}X@{}}
    \toprule
    \multicolumn{1}{c}{\multirow{2}{*}{\textit{Result}}}
      & \multicolumn{3}{c}{\textit{Oracle access}}
      & \multicolumn{1}{c}{\multirow{2}{*}{\textit{Correctness}}}
      & \multicolumn{1}{c}{\multirow{2}{*}{\textit{Status}}} \\
    \cmidrule(lr){2-4}
      & \(\Gen\) & \(\Enc\) & \(\Dec\) & & \\
    \midrule
    \cite{AustrinEtAl22}
      & Q & Q & Q & perfect & Conjectural \\
    \cite{BouazizEtAl24}
      & Q & Q & -- & perfect & Conjectural \\
    \cite{LiEtAl24}
      & C & Q & Q & perfect & Unconditional \\
    \cite{LiEtAl24}
      & C & Q & -- & imperfect & Unconditional \\
    \cite{LiEtAl25}
      & Q (log.) & Q & Q & perfect
      & Unconditional for \(O(\log\lambda)\) \(\Gen\) queries; conjectural
        otherwise \\
    \cite{LiEtAl26PKE}
      & Q & Q & Q & perfect & Unconditional \\
    This work
      & C & Q & Q & imperfect & Unconditional \\
    \bottomrule
  \end{tabularx}
  \caption{Comparison of QPKE impossibility results in the QROM.  Here C and
  Q denote classical and quantum random-oracle access, respectively, and
  -- denotes oracle-free decryption.  The two rows for
  \cite{LiEtAl24} separate its perfect-correctness result from its
  imperfect-correctness result, which requires oracle-free decryption.}
  \label{tab:qpke-comparison}
\end{table}

Moreover, the multi-round theorem gives an interactive key-generation variant of QPKE.  Namely, suppose that ordinary key generation is
replaced by an arbitrary classical-communication setup phase between the
decrypting and encrypting parties.  The parties may perform quantum local
computation and retain quantum states, but both make only classical
random-oracle queries during this phase.  Viewing the setup as the classical
prefix and the challenge ciphertext as Bob's final message, the multi-round
attack applies without any dependence on the number of setup rounds. The reduction fixes two distinct messages and turns the scheme into key
agreement in the natural way. Alice generates and sends the public key; Bob
encrypts one of the two messages and takes its index as his key; Alice
decrypts and recovers the index. Recovering Bob's key is therefore equivalent
to guessing which message was encrypted.

Finally, we obtain the following corollary.

\begin{corollary}[Two-round OSP consequence, informal]
The PKE obtained from a negligibly incorrect two-round OSP protocol by the
transformation of Bartusek and Khurana is not query-bounded IND-CPA secure in
the QROM, provided that the classical OSP sender makes only classical
random-oracle queries and the honest query bounds are polynomial.
\end{corollary}

\subsection{Discussion}

We note that our key-agreement model and security notion are different from prior work.
In classical random-oracle key agreement, a passive eavesdropper can copy the
public transcript and is commonly required to be unable to recover an honest
party's key~\cite{BarakMahmoodyFull}. Here Bob's final message may be quantum
and therefore cannot in general be copied. Moreover, unlike the Classical
Communication One Quantum Message (CC1QM) security notion of Bouaziz-Ermann
et al.~\cite[Definition~3.1]{BouazizEtAl24}, our success condition asks only
that the attacker recover Bob's key. It does not require the attacker to
provide Alice with a replacement state that makes her output the same
non-abort key. Thus our attack objective is weaker, while security against
our interception key-recovery experiment is correspondingly stronger: every
successful CC1QM attack yields a successful interception key-recovery attack,
but the converse need not hold. This does not by itself break the full CC1QM
condition for a quantum final message. It is nevertheless sufficient for our
QPKE application, because the IND-CPA challenge ciphertext is transferred to
the attacker and need not be returned to Alice.

Moreover, our result crucially relies on the fact that the early oracle queries are
classical. Once the complete classical query-answer transcript is fixed, the
party's residual state depends only on that transcript, while all unqueried
oracle values remain uniform. This makes it possible to resample a conditioned
view and then apply random reprogramming. With quantum queries, a party
may instead be entangled with many oracle values in superposition,
and there is no classical transcript that both determines the
residual state and leaves the remaining oracle values independent and
uniform. Measuring the query register generally changes the execution, while
keeping it coherent does not yield the classically sampled
reprogramming set needed by \Cref{lem:reprogramming}. Extending our result to
quantum queries in the prefix would therefore require novel techniques for handling quantum oracle interactions, and we leave this as future work.


\subsection{Acknowledgements}

This research was conducted while AV was a summer intern at NTT Social Informatics Laboratories
in Tokyo. AV is very grateful to NTT and the organizers of the internship for the opportunity and their generous hospitality. 

We acknowledge the use of AI in the production of this paper.  The proof ideas for the two-round setting were human-generated, while GPT-5.6 Sol Extra-Hard Thinking assisted in formalizing. When attempting to generalize to the multi-round setting, GPT-5.6 Sol Extra-Hard Thinking formalized an attack by extending Barak-Mahmoody. The authors take full responsibility for the correctness of this work.

\subsection{Related Work}

\paragraph{Oblivious state preparation.}
Bartusek and Khurana introduced oblivious state preparation (OSP). They showed
that two-round OSP implies one-bit CPA-secure PKE with classical key
generation, decryption, keys, and ciphertexts, but quantum
encryption~\cite[Theorem~7.9]{BartusekKhurana24}. They also reduced OSP to
perfectly complete key agreement, thereby ruling out perfectly correct OSP
in the QROM~\cite[Corollary~7.7]{BartusekKhurana24}. In their relativizing
OSP-to-PKE transformation, the classical OSP sender becomes the key
generator. Our theorem therefore attacks the resulting PKE even when its
correctness error is negligible rather than zero, provided that the sender
queries the random oracle classically.

\paragraph{Quantum public-key encryption.}
Under the Polynomial Compatibility Conjecture, Bouaziz-Ermann et al. handled
protocols with a final quantum message. Their result assumes perfect
completeness and no oracle queries by Alice after she receives that message,
so its QPKE application requires oracle-free
decryption~\cite{BouazizEtAl24}. Our two-message theorem instead permits
imperfect completeness and quantum oracle queries in Alice's final
computation, at the cost of requiring classical queries in her first stage.

Li et al. proved a complementary sequence of QPKE impossibility results.
Their first paper unconditionally rules out perfectly complete QPKE with
classical keys and ciphertexts when key generation has classical oracle
access. It separately rules out imperfectly correct, classical-key QPKE with
possibly quantum ciphertexts in any oracle model, provided that decryption is
oracle-free~\cite{LiEtAl24}. Their next paper allows logarithmically many
quantum key-generation queries unconditionally and gives a conjectural
result beyond that range~\cite{LiEtAl25}. Their 2026 result removes both the
conjecture and the key-generation restriction for perfectly complete QPKE
with classical keys and either classical or quantum
ciphertexts~\cite{LiEtAl26PKE}. Our result keeps the classical-query
restriction on key generation, but allows imperfect correctness and
unrestricted quantum oracle queries by encryption and decryption.

\paragraph{Many-round key agreement.}
In a recent preprint, Li et al. give an unconditional attack against
perfectly complete, arbitrary-round QCCC key agreement with unrestricted
quantum oracle queries~\cite{LiEtAl26ManyRound}. All communication in their
model is classical, and their eavesdropper makes only classical queries. Our
multi-round theorem makes a different tradeoff: it tolerates imperfect
completeness and a final quantum message, but it requires classical oracle
access during the prefix and uses a quantum attacker after the prefix.

\subsection{Organization}

\Cref{sec:technical-overview} explains the two attacks and the QPKE
application.  \Cref{sec:preliminaries} defines the query model, partial
oracles, and the key-agreement and QPKE security notions, and recalls the
external lemmas used in the proofs.  \Cref{sec:resampling} constructs the
heavy-query learner for the two-message setting and defines Alice's
conditioned first-round view.  \Cref{sec:main-attack} proves the two-round and multi-round theorems.  Finally, \Cref{sec:corollaries} derives the classical-eavesdropping,
QPKE, and OSP consequences.

\section{Technical Overview}\label{sec:technical-overview}

\subsection{The Two-Message Attack and QPKE Application}

We consider two-round key agreement in the QROM with the following syntax. Both Alice and Bob are quantum algorithms but Alice may make only classical queries to the oracle and her message to Bob must be classical. Alice first makes at most $q_{\cA, 1}$ queries to the random oracle $H$ to obtain query, answer pairs $T_\cA$, and sends her message $m_1$ to Bob, retaining a residual state $S_\cA$. Bob then makes his queries to $H$, outputs a key $K_B$ and then sends his possibly quantum message $\psi$ to Alice, who then outputs a key $K_A$. 

Eve proceeds as follows. After receiving Alice's first message \(m_1\), Eve looks for an input that is
likely to occur in the transcript of Alice's first round, conditioned on
\(m_1\) and all oracle answers Eve has learned so far. Whenever this
conditional probability exceeds \(\delta\), Eve queries the input and adds
the answer to a partial oracle \(h\), where $h$ is defined such that it agrees with $H$ on all points in its domain, which we denote by $\Dom(h)$. We stop the procedure after
\(\ceil{q_{\cA,1}/\delta^2}\) queries; the probability of reaching this
cutoff is at most \(\delta\). If this learning phase is successful, then every point outside \(\Dom(h)\) 
appears in a fresh conditioned Alice transcript with probability at most
\(\delta\). 

Although Alice's queries in the first round are classical, she may retain a quantum
state. However, once Alice's message and full query transcript
are fixed, her residual state depends only on those classical random variables, while
all unqueried oracle values remain uniform. Conditioned on \((m_1,h)\), Eve
can therefore prepare a distribution over Alice's views, which we denote by $\omega_{m_1, h}$, and use it to sample a new view
\[
  (T'_{\cA},S'_{\cA})\leftarrow\omega_{m_1,h}.
\] 
Eve then intercepts Bob's response \(\Psi\) and runs Alice's final algorithm
on \(S'_{\cA}\) and \(\Psi\), using a \emph{reprogrammed oracle}
\(H[T'_{\cA}]\), which is defined to agree with $T'_\cA$ on all points inside of $\Dom(T'_{\cA})$, and with $H$ everywhere else. 

The analysis of the two-message attack combines two ingredients. The
first ingredient is the heavy-query bound of Austrin et
al.~\cite[Lemma~3.4]{AustrinEtAl22}, in the form recalled by Katz and
Sela~\cite[Lemma~4]{KatzSela24}, which lets Eve learn the heavy inputs
before we invoke reprogramming. The second is the random
reprogramming lemma in the form stated by Katz and
Sela~\cite[Lemma~3]{KatzSela24}. The lemma bounds the effect
of reprogramming an oracle on a random set: if each point belongs to that set
with probability at most \(\delta\), then the acceptance probability of a
\(q\)-query quantum algorithm changes by at most \(2q\sqrt\delta\). 

The QPKE application uses the standard reduction from QPKE to key
agreement used in prior QROM impossibility work
\cite{BouazizEtAl24,LiEtAl24}. Alice generates
\((\pk,\sk)\) and sends \(\pk\). Fix two distinct messages \(m_0,m_1\).
Bob chooses a uniform bit \(b\), encrypts \(m_b\), and takes \(b\) as his
key; Alice decrypts and maps \(m_0,m_1\) back to their indices. The resulting
agreement probability is at least the scheme's minimum correct-decryption
probability. In the IND-CPA game, Eve preprocesses \(\pk\), treats the
challenge ciphertext as Bob's response, and uses its recovered key as the
guess for \(b\). Interception success \(p\) therefore yields IND-CPA
advantage at least \(p-1/2\). Because the ciphertext is transferred rather
than copied, encryption, decryption, and the ciphertext may all be quantum.
Only key generation needs classical oracle access.

\subsection{The Multi-Round Attack}

For the multi-round result, we treat the entire classical-query prefix as an
auxiliary classical protocol.  The parties' local computations and residual
states may be quantum, but their messages and oracle queries during the
prefix are classical.  Since we do not impose a computational bound, an
exact simulator can sample these classical actions from their induced
distributions and record the corresponding conditional residual states.  We
can therefore apply the Barak--Mahmoody learner to the simulated classical
views and then map its guarantees back to the parties' actual query
transcripts and residual quantum states~\cite{BarakMahmoodyFull}.

Conditioned on the public transcript \(m\) and the partial oracle \(h\)
learned by Eve, the parties' prefix views are close to the product of their
marginals, and no input outside \(\Dom(h)\) is likely to appear in Alice's
prefix query transcript.  This product guarantee does not yet include the
random oracle used in the final round.  Since every prefix query is
classical, after the two query transcripts are fixed, the unqueried oracle
values remain uniform and the residual states have no further dependence on
them.  We use this fact to sample independent prefix views together with a
consistent oracle and show that the resulting state is close to the real
conditioned state.  These steps are given in
\Cref{sec:classical-prefix-views,sec:bm-learner,sec:conditioned-prefix-resampling}.

Eve can then sample a fresh Alice prefix view, intercept Bob's final message,
and run Alice's final computation from the sampled residual state.  Since
the learner's query bound depends only on the total numbers of prefix
queries made by Alice and Bob, the attack has no additional dependence on
the number of prefix rounds.


\section{Preliminaries}\label{sec:preliminaries}

Let \(\lambda\in\mathbb{N}\) be the security parameter, and let
\(n,\ell:\mathbb{N}\to\mathbb{N}_{>0}\) be functions bounded above by a
polynomial in \(\lambda\). Define
\[
  \cX_\lambda=\bits^{n(\lambda)},
  \qquad
  \cY_\lambda=\bits^{\ell(\lambda)},
  \qquad
  \cH_\lambda=\{H:\cX_\lambda\to\cY_\lambda\}.
\]
For any finite set \(\cR\), we write \(r\leftarrow\cR\) for a uniformly
random element of \(\cR\).  In particular,
\(H\leftarrow\cH_\lambda\) denotes a uniformly random function.  When
\(\lambda\) is clear, we omit the subscript.  When representing the
sampled oracle by a classical register, \(\mathsf H\) has orthonormal basis
\(\{\ket{H}_{\mathsf H}:H\in\cH\}\). For a probability mass function or subprobability mass function \(\pi\), we
write
\[
  \Supp(\pi):=\{z:\pi(z)>0\}
\]
for its support.

\subsection{Random Oracles and Query Algorithms}

\begin{definition}[Oracle quantum algorithms]\label{def:oracle-algorithms}
A quantum oracle algorithm
\(\mathcal Q=\{\mathcal Q_\lambda\}_{\lambda\in\mathbb N}\) is a family
of quantum circuits. We write \(I\) for its input register, \(X\) for its
\(n(\lambda)\)-qubit query register, \(Y\) for its
\(\ell(\lambda)\)-qubit response register, and \(Z\) for all remaining work
registers. The circuit family depends only on $\lambda$. For
\(H\in\cH_\lambda\), a query applies the unitary
\[
  U_H\ket{x}_X\ket{y}_Y\ket{z}_Z
  =
  \ket{x}_X\ket{y\oplus H(x)}_Y\ket{z}_Z.
\]
The algorithm alternates oracle-independent operations with applications of
\(U_H\). For fixed \(\lambda\), let \(q(\mathcal Q_\lambda)\) denote its
worst-case number of oracle queries; we sometimes abbreviate this to
\(q(\mathcal Q)\). 

An algorithm has \emph{classical oracle access} if each query is made on a
classical input \(x\), the answer \(H(x)\) is returned classically, and the
ordered list of query-answer pairs can be recorded without changing the
execution. The rest of the computation and the residual state may be
quantum.
\end{definition}

We use the standard classical-quantum representation of a state that is
classical on one register. If \(C\)
is a classical register with finite alphabet \(\mathcal C\) and \(S\) is a
quantum register, then
\[
  \rho_{CS}
  =
  \sum_{c\in\mathcal C}
  p(c)\proj{c}_C\otimes\sigma_{S\mid c},
\]
where \(p\) is a probability distribution and, whenever \(p(c)>0\),
\(\sigma_{S\mid c}\) is the normalized state of \(S\) obtained by measuring
\(C\) in the computational basis and observing \(c\). Explicitly, writing
\(P_c=\proj{c}_C\),
\[
  p(c)=\Tr\!\left[(P_c\otimes I_S)\rho_{CS}\right]
\]
and, for \(p(c)>0\),
\[
  \sigma_{S\mid c}
  =
  \frac{
    \Tr_C\!\left[(P_c\otimes I_S)\rho_{CS}(P_c\otimes I_S)\right]
  }{p(c)}.
\]

\subsection{Partial Oracles and Reprogramming}

\begin{definition}[Partial oracle]
For any $\lambda \in \mathbb{N}$, fix $\cX = \cX_\lambda, \cY = \cY_\lambda,$ and $\cH = \cH_\lambda$.  A partial oracle is a function
\(h:\Dom(h)\to\cY\), where
\(\Dom(h)\subseteq\cX\). We denote the \emph{size} of $h$ by \(|h|=|\Dom(h)|\). We write
\(h\preceq H\) when a full oracle \(H\in\cH\) extends \(h\), meaning that
\(H(x)=h(x)\) for every \(x\in\Dom(h)\).  Define
\[
  \cH(h)=\{H\in\cH:h\preceq H\}.
\]
Two partial oracles \(h\) and \(\tau\) are \emph{consistent} if they agree on the
intersection of their domains.  For consistent \(h\) and \(\tau\), define
\[
  \cH(h,\tau)
  =
  \{H\in\cH:h\preceq H,\ \tau\preceq H\}.
\]
The notation \(H\leftarrow\Unif(\cH(h))\), and similarly
\(H\leftarrow\Unif(\cH(h,\tau))\), means that \(H\) is sampled uniformly
from the indicated set.
\end{definition}

An ordered transcript
\(\tau=((x_1,y_1),\ldots,(x_t,y_t))\) is \emph{consistent} when
\(x_i=x_j\) implies \(y_i=y_j\).  A consistent transcript defines the
partial oracle \(\tau(x_i)=y_i\).  Repeated queries may be retained in the
ordered list, while \(\Dom(\tau)\) denotes the set of distinct queried
points.  We write \(|\tau|=t\) for the length of the ordered list, including
repetitions.

\begin{definition}[Reprogrammed oracle]
For an oracle \(H\) and a consistent transcript \(\tau\) (equivalently, a
partial oracle), define
\[
  H[\tau](x)
  =
  \begin{cases}
    \tau(x),&x\in\Dom(\tau),\\
    H(x),&x\notin\Dom(\tau).
  \end{cases}
\]
For another internally consistent transcript \(\tau'\), which need not be
consistent with \(\tau\), successive reprogramming is
defined so that the second transcript overwrites the first on their common
domain:
\[
  H[\tau][\tau'](x)
  =
  \begin{cases}
    \tau'(x),&x\in\Dom(\tau'),\\
    \tau(x),&x\in\Dom(\tau)\setminus\Dom(\tau'),\\
    H(x),&x\notin\Dom(\tau)\cup\Dom(\tau').
  \end{cases}
\]
\end{definition}


\subsection{External Lemmas}

Below we recall the two prior results used in our reprogramming attack.

The first is the random reprogramming
lemma in the form used by Katz and Sela~\cite[Lemma~3]{KatzSela24}, itself a
restatement of a result by Alagic et al.~\cite[Lemma~3]{AlagicEtAl22}.

\begin{lemma}[Random reprogramming lemma]\label{lem:reprogramming}
Let \(\cD\) be a distinguisher in the following game.
\begin{enumerate}[leftmargin=2em]
  \item \(\cD\) outputs a full function
  \(F_0:\cX\to\cY\) and a possibly inefficient randomized classical
  algorithm \(C\) whose output is a
  partial function \(R\).  Let
  \[
    \varepsilon
    =
    \max_{x\in\cX}
    \Pr_{R\leftarrow C}[x\in\Dom(R)].
  \]
  \item The challenger samples \(R\leftarrow C\), sets
  \(F_1=F_0[R]\), chooses \(b\leftarrow\bits\), and gives \(\cD\) quantum
  oracle access to \(F_b\).
  \item The oracle access is removed.  The challenger reveals the random
  coins used by \(C\), and \(\cD\) outputs a bit.
\end{enumerate}
If \(\cD\) makes at most \(q\) challenge-oracle queries, then
\[
  \left|
    \Pr[\cD\text{ outputs }1\mid b=1]
    -
    \Pr[\cD\text{ outputs }1\mid b=0]
  \right|
  \le 2q\sqrt\varepsilon.
\]
\end{lemma}

The restriction that \(C\) is classical is deliberate.  Its only role is to
sample the classical partial function \(R\), and the game later reveals the
random coins used for that sample.  In our applications, \(C\) samples a
classical query transcript, while \(\cD\) separately prepares any residual
quantum state conditioned on the revealed transcript.  Thus \(C\) never
simulates Alice's quantum state.  Allowing \(C\) itself to retain quantum
state would also require a different game interface, since revealing
classical random coins would no longer specify its complete internal state.

The second is the following heavy-query counting lemma of Austrin et
al.~\cite[Lemma~3.4]{AustrinEtAl22}, in the form recalled by Katz and
Sela~\cite[Lemma~4]{KatzSela24}.

\begin{lemma}[Heavy-query bound]\label{lem:heavy-bound}
Let \(L,z_1,x_1,\ldots,z_q,x_q\) be correlated random variables, where
\(L\) ranges over subsets of a universe \(\cU\), each \(x_i\in\cU\), and
the \(x_i\)'s are pairwise distinct.  Conditioned on a history
\(z_1,x_1,\ldots,z_i\), call \(x_i\) \(\delta\)-heavy if
\[
  \Pr[x_i\in L\mid z_1,x_1,\ldots,z_i]\ge\delta.
\]
Let \(S\) be the set of \(\delta\)-heavy points among the \(x_i\)'s.  Then
\[
  \mathbb{E}[|S|]\le\frac{\mathbb{E}[|L|]}{\delta}.
\]
\end{lemma}

\subsection{Public Key Cryptography with Quantum Parties}\label{sec:model}

We now define the protocols used throughout
the paper. We note that the two-message model in \Cref{def:protocol} is close to the CC1QM
model of Bouaziz-Ermann et al.~\cite{BouazizEtAl24}, with one additional
restriction: before sending her first message, Alice has only classical
oracle access.

\subsubsection{Two-Message Key Agreement}


\begin{definition}[Two-message key agreement with a classical first message]
\label{def:protocol}
A two-message key-agreement protocol with a classical first message and a
possibly quantum response is a tuple of oracle quantum algorithms \(\Pi=(\cA_1,\cB,\cA_2)\) with the
following syntax relative to \(\cH_\lambda\). 
\begin{itemize}[leftmargin=2em]
  \item \(\cA_1^H(1^\lambda)\) has classical oracle access and outputs
  \((T_{\cA},M_1,S_{\cA})\).  Here \(M_1\) is a classical first message,
  \(S_{\cA}\) is Alice's residual register, and \(T_{\cA}\) is the complete
  ordered classical query-answer transcript of \(\cA_1\).
  \item \(\cB^H(1^\lambda,M_1)\) has quantum oracle access and outputs a
  possibly quantum message register \(\Psi\) and a classical key
  \(K_{\cB}\in\cK_\lambda\cup\{\bot\}\).
  \item \(\cA_2^H(1^\lambda,S_{\cA},\Psi)\) has quantum oracle access and
  outputs a classical key \(K_{\cA}\in\cK_\lambda\cup\{\bot\}\).
\end{itemize}
Other than access to the common random oracle, Alice and Bob initially share
no state or correlated randomness.  We assume worst-case bounds
\[
  |T_{\cA}|\le q_{\cA,1},
  \qquad
  q(\cB)\le q_{\cB},
  \qquad
  q(\cA_2)\le q_{\cA,2}.
\] 
The protocol is the following experiment, denoted by \(\Real_\Pi(1^\lambda)\):
\begin{enumerate}[leftmargin=2em]
  \item Sample \(H\leftarrow\cH_\lambda\).
  \item Run
  \((T_{\cA},M_1,S_{\cA})\leftarrow\cA_1^H(1^\lambda)\), and send the
  classical message \(M_1\) to Bob.
  \item Run
  \((\Psi,K_{\cB})\leftarrow\cB^H(1^\lambda,M_1)\), and send the
  message register \(\Psi\) to Alice.
  \item Run
  \(K_{\cA}\leftarrow\cA_2^H(1^\lambda,S_{\cA},\Psi)\).
\end{enumerate}
The valid agreement probability is
\[
  \alpha_\Pi(\lambda)
  :=
  \Pr_{\Real_\Pi(1^\lambda)}
  [K_{\cA}=K_{\cB}\neq\bot].
\]
We use the following notions of completeness and security for such a protocol:

\begin{enumerate}[leftmargin=2em,label=\textup{(\roman*)}]
  \item \textbf{Completeness.}  The protocol has completeness error
  \(\epsilon_c(\lambda)\) if
  \[
    \alpha_\Pi(\lambda)\ge 1-\epsilon_c(\lambda).
  \]
  We say the protocol is \emph{correct} if \(\epsilon_c\) is negligible.
  \item \textbf{Security.}  The protocol is interception key-recovery secure,
  as defined in \Cref{def:intkr}, if, for every polynomially query-bounded attacker
  \(\cE\),
  \begin{equation}
    \Adv^{\IntKR}_{\Pi,\cE}(\lambda) \leq \negl(\lambda).
  \end{equation}

\end{enumerate}
\end{definition}

\begin{definition}[Intercepting attacker]
An \emph{intercepting-attacker} is a pair
\(\cE=(\cE_{1},\cE_{2})\). Its first stage receives
Alice's classical message, may query the oracle before Bob generates his
response, and outputs a possibly quantum state \(Z_{\cE}\). Its second stage
receives \(Z_{\cE}\) together with Bob's message register and outputs a guess for the
classical key.
\end{definition}

\begin{definition}[Interception key-recovery experiment]\label{def:intkr}
Let \(\cE=(\cE_{1},\cE_{2})\) be any polynomially
query-bounded intercepting-attacker: its computation may be
unbounded, but its two stages together make at most \(\poly(\lambda)\)
random-oracle queries.  The experiment
\(\IntKR_{\Pi,\cE}(1^\lambda)\) is:
\begin{enumerate}[leftmargin=2em]
  \item Sample \(H\leftarrow\cH_\lambda\).
  \item Run
  \((T_{\cA},M_1,S_{\cA})\leftarrow\cA_1^H(1^\lambda)\) and reveal
  \(M_1\) to \(\cE_{1}\).
  \item Run
  \(Z_{\cE}\leftarrow\cE_{1}^H(1^\lambda,M_1)\).
  \item Run
  \((\Psi,K_{\cB})\leftarrow\cB^H(1^\lambda,M_1)\) and transfer
  \(\Psi\) to \(\cE_{2}\).
  \item Run
  \(K^*\leftarrow\cE_{2}^H(1^\lambda,Z_{\cE},\Psi)\).
  The experiment outputs \(1\) iff \(K^*=K_{\cB}\neq\bot\).
\end{enumerate}
Define
\[
  \Adv^{\IntKR}_{\Pi,\cE}(\lambda)
  :=
  \Pr[\IntKR_{\Pi,\cE}(1^\lambda)=1].
\]
\end{definition}

\subsubsection{Multi-Round Key Agreement with a Classical-Query Prefix}

We also consider protocols with an arbitrary classical-communication prefix.
During the prefix, both parties query the oracle only
classically, although their other computations and residual states may be
quantum. In the final round, Bob sends a possibly quantum message, after
which Alice performs her final computation. This communication pattern is similar to
the CC1QM pattern of Bouaziz-Ermann et
al.~\cite[Definition~11]{BouazizEtAl24}, except that we additionally require classical
oracle access during the prefix.

\begin{definition}[Multi-round key agreement with a classical-query prefix]\label{def:multi-protocol}

A key-agreement protocol with a classical prefix and one final possibly
quantum message is a tuple of oracle quantum algorithms
\[
  \Pi^{\mathsf{mr}}
  =
  \bigl((\cA_i)_{i=1}^{r},(\cB_i)_{i=1}^{r-1},
    \cB_{\mathsf q},\cA_{\mathsf{fin}}\bigr).
\]
Here \(r\ge1\) is the number of messages that Alice sends during the
classical prefix. To simplify the indexing, we assume that Alice sends both
the first and last prefix messages.
The honest experiment, denoted
\(\Real_{\Pi^{\mathsf{mr}}}(1^\lambda)\), is:
\begin{enumerate}[leftmargin=2em]
  \item Sample \(H\leftarrow\cH_\lambda\). Alice and Bob begin in an
  oracle-independent product state
  \(\rho_{\cA,0}\otimes\rho_{\cB,0}\).
  \item For \(i=1,\ldots,r\), Alice runs \(\cA_i\) on her current residual
  register and the public transcript so far, makes only classical oracle
  queries, sends a classical message \(M_{2i-1}\), and retains her new
  residual register. If \(i<r\), Bob proceeds analogously: he runs
  \(\cB_i\), makes only classical oracle queries, sends a classical message
  \(M_{2i}\), and retains his new residual register.
  \item Let
  \[
    M=(M_1,\ldots,M_{2r-1})
  \]
  be the complete classical prefix transcript. Let \(T_{\cA}\) and
  \(T_{\cB}\) be the complete ordered query-answer transcripts accumulated by
  Alice and Bob during the prefix, and let \(S_{\cA}\) and \(S_{\cB}\) be
  their residual quantum registers at its end.
  \item Bob runs
  \[
    (\Psi,K_{\cB})
    \leftarrow
    \cB_{\mathsf q}^{H}
    (1^\lambda,M,S_{\cB}),
  \]
  where \(\Psi\) is a possibly quantum message register and
  \(K_{\cB}\in\cK_\lambda\cup\{\bot\}\) is classical.
  \item After receiving \(\Psi\), Alice runs
  \[
    K_{\cA}
    \leftarrow
    \cA_{\mathsf{fin}}^{H}
    (1^\lambda,M,S_{\cA},\Psi).
  \]
  Here \(K_{\cA}\in\cK_\lambda\cup\{\bot\}\) is classical.
\end{enumerate}
Write
\[
  q_{\cA,\mathsf{pre}}
  :=\max |T_{\cA}|,
  \qquad
  q_{\cB,\mathsf{pre}}
  :=\max |T_{\cB}|,
\]
and let \(q_{\cB,\mathsf q}\) and
\(q_{\cA,\mathsf{fin}}\) bound the quantum-oracle queries made by
\(\cB_{\mathsf q}\) and \(\cA_{\mathsf{fin}}\), respectively. The valid
agreement probability is
\[
  \alpha_{\Pi^{\mathsf{mr}}}(\lambda)
  :=
  \Pr_{\Real_{\Pi^{\mathsf{mr}}}(1^\lambda)}
  [K_{\cA}=K_{\cB}\neq\bot].
\]
We use the following notions of completeness and security for such a protocol:
\begin{enumerate}[leftmargin=2em,label=\textup{(\roman*)}]
  \item \textbf{Completeness.}  The protocol has completeness error
  \(\epsilon_c(\lambda)\) if
  \[
    \alpha_{\Pi^{\mathsf{mr}}}(\lambda)
    \ge 1-\epsilon_c(\lambda).
  \]
  We say the protocol is \emph{correct} if \(\epsilon_c\) is negligible.
  \item \textbf{Security.}  The protocol is online interception key-recovery
  secure, as defined in \Cref{def:multi-intkr}, if, for every
  computationally unbounded attacker \(\cE\) making at most \(\poly(\lambda)\)
  random-oracle queries, restricted to classical queries during the prefix
  but allowed quantum queries after receiving \(\Psi\),
  \begin{equation}
    \Adv^{\IntKR,\mathsf{mr}}_{\Pi^{\mathsf{mr}},\cE}(\lambda)
    \leq \negl(\lambda).
  \end{equation}
\end{enumerate}
\end{definition}

\begin{definition}[Online interception experiment]
\label{def:multi-intkr}
The experiment
\(\IntKR^{\mathsf{mr}}_{\Pi^{\mathsf{mr}},\cE}(1^\lambda)\) proceeds as
follows.
\begin{enumerate}[leftmargin=2em]
  \item Sample \(H\leftarrow\cH_\lambda\) and initialize the parties and Eve.
  \item Run the honest classical prefix. After each public message, activate
  Eve on the transcript so far and its retained state; during these
  activations Eve may make classical oracle queries and update its state.
  \item Run \(\cB_{\mathsf q}^H\) on Bob's residual register. Transfer its
  message register \(\Psi\) to Eve and retain Bob's key \(K_{\cB}\).
  \item Eve may make quantum oracle queries and outputs a classical key guess
  \(K^*\). The experiment outputs \(1\) iff
  \(K^*=K_{\cB}\neq\bot\).
\end{enumerate}
Define
\[
  \Adv^{\IntKR,\mathsf{mr}}_{\Pi^{\mathsf{mr}},\cE}(\lambda)
  :=
  \Pr[\IntKR^{\mathsf{mr}}_{\Pi^{\mathsf{mr}},\cE}(1^\lambda)=1].
\]
\end{definition}

\subsubsection{Quantum Public-Key Encryption}

We consider classical public and secret keys, classical messages whose
length is bounded by a polynomial in \(\lambda\), and possibly quantum encryption,
decryption, and ciphertexts. We allow imperfect
correctness but require \(\Gen\) to query the random oracle only classically.
Our security experiment is the corresponding query-bounded IND-CPA
experiment: the attacker may perform unbounded computation but is charged
for its oracle queries. When \(\Gen\), \(\Dec\), and the ciphertext are
classical but \(\Enc\) may be quantum, the \(\mu(\lambda)=1\) case
is exactly the PKE syntax used by Bartusek and
Khurana~\cite[Definition~7.8]{BartusekKhurana24}.

\begin{definition}[QPKE with classical-query key generation]
\label{def:qpke}
Let \(\mu:\mathbb{N}\to\mathbb{N}\) be any polynomial-time computable
function satisfying \(1\le\mu(\lambda)\le p(\lambda)\) for some polynomial
\(p\), and let
\(\cM_\lambda=\bits^{\mu(\lambda)}\). A quantum public-key encryption scheme
\(\Sigma=(\Gen,\Enc,\Dec)\) for the message spaces
\(\{\cM_\lambda\}_{\lambda\in\mathbb N}\), in the QROM with oracle family
\(\cH_\lambda\), has the following syntax.
\begin{itemize}[leftmargin=2em]
  \item \((\pk,\sk)\leftarrow\Gen^H(1^\lambda)\) outputs a classical
  public key and secret key.  The algorithm has classical random-oracle
  access, and its complete ordered query-answer transcript can be recorded as
  \(T_{\Gen}\).
  \item \(C\leftarrow\Enc^H(1^\lambda,\pk,m)\), for
  \(m\in\cM_\lambda\), uses quantum random-oracle access and outputs a
  possibly quantum ciphertext register \(C\).
  \item \(\widehat m\leftarrow\Dec^H(1^\lambda,\sk,C)\) uses quantum
  random-oracle access and outputs
  \(\widehat m\in\cM_\lambda\cup\{\bot\}\).
\end{itemize}
Write \(q_{\Gen}\), \(q_{\Enc}\), and \(q_{\Dec}\) for the respective
worst-case query bounds. Define the minimum correct-decryption probability
\begin{align*}
  \alpha_\Sigma(\lambda)
  :=\min_{m\in\cM_\lambda}
  \Pr\bigl[&\Dec^H(1^\lambda,\sk,C)=m:\ \\
    &H\leftarrow\cH_\lambda,\
    (\pk,\sk)\leftarrow\Gen^H(1^\lambda),\
    C\leftarrow\Enc^H(1^\lambda,\pk,m)\bigr].
\end{align*}
The scheme has correctness error \(\epsilon_c(\lambda)\) if
\(\alpha_\Sigma(\lambda)\ge1-\epsilon_c(\lambda)\).
\end{definition}

\begin{definition}[Query-bounded IND-CPA security]\label{def:qpke-ind}
For a two-stage attacker \(\cD=(\cD_0,\cD_1)\), the experiment
\(\IND_{\Sigma,\cD}(1^\lambda)\) is:
\begin{enumerate}[leftmargin=2em]
  \item Sample \(H\leftarrow\cH_\lambda\) and
  \((\pk,\sk)\leftarrow\Gen^H(1^\lambda)\).
  \item Run
  \((Z_{\cD},m_0,m_1)\leftarrow\cD_0^H(1^\lambda,\pk)\), where
  \(Z_{\cD}\) is a possibly quantum state and
  \(m_0,m_1\in\cM_\lambda\).
  \item Sample \(b\leftarrow\bits\), run
  \(C\leftarrow\Enc^H(1^\lambda,\pk,m_b)\), and transfer \(C\) to
  \(\cD_1\).
  \item Run \(b'\leftarrow\cD_1^H(1^\lambda,Z_{\cD},C)\).  The experiment
  outputs \(1\) iff \(b'=b\).
\end{enumerate}
Define
\[
  \Adv^{\IND\text{-}\mathrm{CPA}}_{\Sigma,\cD}(\lambda)
  :=
  \left|\Pr[\IND_{\Sigma,\cD}(1^\lambda)=1]-\frac12\right|.
\]
The scheme is query-bounded IND-CPA secure if this advantage is negligible
for every computationally unbounded attacker making at most
\(\poly(\lambda)\) random-oracle queries. This differs from standard
computational IND-CPA only in the attacker's resource bound. Because
encryption is public, the experiment does not need a separate encryption
oracle.
\end{definition}

\section{Conditional Resampling and Heavy-Query Learning}\label{sec:resampling}

\subsection{The Heavy-Query Attacker}

The following procedure is the \textsc{findTranscript} heavy-query attacker
of Katz and Sela~\cite[Section~3 and Lemma~5]{KatzSela24} adapted to our setting. Its cutoff analysis is precisely
their application of the heavy-query counting bound of Austrin et
al.~\cite{AustrinEtAl22}.

Throughout this subsection, probabilities are taken over
\[
  H\leftarrow\cH,
  \qquad
  (T_{\cA},M_1,S_{\cA})\leftarrow\cA_1^H(1^\lambda).
\]
For a first-message value \(m_1\) and a partial oracle \(h\), let
\[
  \mathsf E_{m_1,h}:=\{M_1=m_1,\ h\preceq H\}.
\]
Whenever \(\Pr[\mathsf E_{m_1,h}]>0\), define
\begin{equation}\label{eq:p-heavy}
  p_{m_1,h}(x)
  :=
  \Pr[x\in\Dom(T_{\cA})\mid\mathsf E_{m_1,h}].
\end{equation}
Define
\[
  \Heavy_\delta(m_1,h)
  :=
  \{x\in\cX\setminus\Dom(h):p_{m_1,h}(x)>\delta\}.
\]
Fix a total order \(\prec\) on \(\cX\).

\begin{algorithm}[H]
\caption{\(\Learn_\delta^H(1^\lambda,m_1)\)}
\label{alg:learn}
\KwInput{Security parameter \(1^\lambda\), first message \(m_1\), oracle
access to \(H\), threshold \(\delta\), and bound \(q_{\cA,1}\).}
\KwOutput{A partial oracle \(h\) or \(\bot\).}
\(N_{\max}\gets\ceil{q_{\cA,1}/\delta^2}\); \(h\gets\emptyset\)\;
\While{\(\Heavy_\delta(m_1,h)\neq\emptyset\)}{
  \If{\(|\Dom(h)|\ge N_{\max}\)}{\Return \(\bot\)\;}
  Let \(x\) be the \(\prec\)-least point in \(\Heavy_\delta(m_1,h)\)\;
  Query \(H(x)\) and set \(h\gets h\cup\{x\mapsto H(x)\}\)\;
}
\Return \(h\)\;
\end{algorithm}

A pair \((m_1,h)\) is \emph{reachable} if
\[
  \Pr[M_1=m_1,\ \Learn_\delta^H(1^\lambda,m_1)=h]>0.
\]

As in the Katz-Sela attacker, fixing the order \(\prec\) makes every next
query a deterministic function of \((m_1,h)\). Hence, for every reachable
pair \((m_1,h)\) and every oracle \(H\),
\[
  \Learn_\delta^H(1^\lambda,m_1)=h
  \quad\Longleftrightarrow\quad
  h\preceq H.
\]
We refer to this equivalence as the \emph{deterministic-output property} of
\(\Learn\).

\subsection{The Conditioned View of the First Round}

For every reachable pair \((m_1,h)\), let \(\omega_{m_1,h}\)
denote the conditional state of Alice's view after the first round
\((T_{\cA},S_{\cA})\) given \(\mathsf E_{m_1,h}\).  Write
\[
  \mu_{m_1,h}(\tau)
  =
  \Pr[T_{\cA}=\tau\mid\mathsf E_{m_1,h}]
\]
for its transcript marginal.  Explicitly,
for every \(\tau\) with \(\mu_{m_1,h}(\tau)>0\), let
\(\sigma_{m_1,\tau}\) be the corresponding conditional state of
\(S_{\cA}\) given \((M_1,T_{\cA})=(m_1,\tau)\). Then
\[
  \omega_{m_1,h}
  =
  \sum_{\substack{\tau:\\
        \mu_{m_1,h}(\tau)>0}}
  \mu_{m_1,h}(\tau)
  \proj{\tau}_{T_{\cA}}\otimes\sigma_{m_1,\tau}.
\]
The notation
\((T'_{\cA},S'_{\cA})\leftarrow\omega_{m_1,h}\) means that we prepare a new
instance of this state: sample \(\tau\leftarrow\mu_{m_1,h}\), place \(\tau\)
in \(T'_{\cA}\), and prepare \(S'_{\cA}\) in the state
\(\sigma_{m_1,\tau}\). By the deterministic-output property, the events
\(\mathsf E_{m_1,h}\) and
\(\{M_1=m_1,\Learn_\delta^H(1^\lambda,m_1)=h\}\) coincide for every
reachable pair.

\begin{lemma}[Katz-Sela heavy-query attacker]\label{lem:learn-success}
For every \(0<\delta<1\), with probability at least \(1-\delta\) over the
real execution of the first round, \(\Learn_\delta^H(1^\lambda,M_1)\) outputs a
partial oracle \(h\neq\bot\).  On every nonabort execution,
\[
  |\Dom(h)|\le\ceil{q_{\cA,1}/\delta^2}
\]
and, for every \(x\notin\Dom(h)\),
\[
  \Pr_{(T'_{\cA},S'_{\cA})\leftarrow\omega_{M_1,h}}
  [x\in\Dom(T'_{\cA})]
  \le\delta.
\]
\end{lemma}

\begin{proof}
The argument is a restatement, in our notation, of the proof of Katz and
Sela~\cite[Lemma~5]{KatzSela24}; we include it explicitly for completeness.
The residual register \(S_{\cA}\) plays no role in the counting argument and
is simply prepared together with its sampled transcript.

The second claim is immediate upon successful termination:
\(\Heavy_\delta(M_1,h)=\emptyset\), and the displayed probability is exactly
\(p_{M_1,h}(x)\).

We now bound the abort probability. If \(q_{\cA,1}=0\), the attacker
returns the empty partial oracle. Assume \(q_{\cA,1}>0\), and let
\(N_{\max}=\ceil{q_{\cA,1}/\delta^2}\). Let
\(L=\Dom(T_{\cA})\). Let \(N\le N_{\max}\) be the number of real queries
issued by the attacker before it returns or aborts. On abort,
\(N=N_{\max}\); successful termination with \(N=N_{\max}\) is also allowed.

To apply \Cref{lem:heavy-bound} with a fixed-length sequence, set
\(z_1=M_1\). For each real attacker step \(i\le N\), let \(x_i\) be its
query and let \(z_{i+1}=H(x_i)\). For \(N<i\le N_{\max}\), pad with pairwise
distinct symbols \(x_i=\bot_i\notin\cX\) and fixed values
\(z_{i+1}=0\). Work over the enlarged universe
\(\cX\sqcup\{\bot_1,\ldots,\bot_{N_{\max}}\}\). The points are pairwise
distinct because the attacker never repeats a query.

Immediately before a real query \(x_i\), the history determines the current
partial oracle \(h_{i-1}\). Hence
\[
  \Pr[x_i\in L\mid z_1,x_1,\ldots,z_i]
  =p_{M_1,h_{i-1}}(x_i)>\delta.
\]
Every real attacker query is therefore \(\delta\)-heavy, while every padding
point has conditional membership probability zero. The heavy set in
\Cref{lem:heavy-bound} has size exactly \(N\), and so
\[
  \mathbb{E}[N]
  \le
  \frac{\mathbb{E}[|L|]}{\delta}
  \le
  \frac{q_{\cA,1}}{\delta}.
\]
By Markov's inequality,
\[
  \Pr[N\ge N_{\max}]
  \le
  \frac{q_{\cA,1}}{\delta N_{\max}}
  \le\delta.
\]
Since abort implies \(N=N_{\max}\),
\[
  \Pr[\Learn_\delta^H(1^\lambda,M_1)=\bot]
  \le \Pr[N\ge N_{\max}]
  \le \delta.
\]
\end{proof}

\section{Reprogramming Attacks}\label{sec:main-attack}

\subsection{The Two-Message attacker}

\begin{definition}[Eve's two-stage reprogramming attack]\label{def:eve}
Fix \(0<\delta<1\).  The attacker
\(\cE_\delta=(\cE_{\delta,1},
\cE_{\delta,2})\) is defined as follows.

On input the first message \(m_1\), the preprocessing stage runs
\[
  h\leftarrow\Learn_\delta^H(1^\lambda,m_1).
\]
If \(h=\bot\), it outputs \(Z_{\cE}=\bot\).  Otherwise it prepares
\[
  (T'_{\cA},S'_{\cA})\leftarrow\omega_{m_1,h}
\]
and outputs \(Z_{\cE}=(m_1,h,T'_{\cA},S'_{\cA})\).

After receiving Bob's message register \(\Psi\), the postprocessing stage
outputs \(\bot\) if \(Z_{\cE}=\bot\).  Otherwise it defines the fixed
reprogrammed oracle \(H'=H[T'_{\cA}]\), runs
\[
  K^*\leftarrow
  \cA_2^{H'}(1^\lambda,S'_{\cA},\Psi),
\]
and outputs \(K^*\).
\end{definition}

The preprocessing stage makes at most
\(\ceil{q_{\cA,1}/\delta^2}\) classical queries. Each query to
\(H[T'_{\cA}]\) can be implemented using one query to \(H\), a reversible
lookup in the known transcript, and a \(\ket{+}^{\otimes\ell}\) ancilla; see
the proof of \Cref{lem:g3-g4}. The postprocessing stage therefore makes at
most \(q_{\cA,2}\) quantum queries.

\begin{theorem}\label{thm:main}
Let \(\Pi=(\cA_1,\cB,\cA_2)\) be a protocol as in
\Cref{def:protocol}, and let
\(\alpha=\alpha_\Pi(\lambda)\).  For every \(0<\delta<1\), the attacker
\(\cE_\delta\) of \Cref{def:eve} satisfies
\[
  \Adv^{\IntKR}_{\Pi,\cE_\delta}(\lambda)
  \ge
  \alpha
  -\delta
  -2(2q_{\cB}+q_{\cA,2})\sqrt\delta.
\]
It makes at most
\[
  \ceil{q_{\cA,1}/\delta^2}+q_{\cA,2}
\]
random-oracle queries.
\end{theorem}

\begin{proof}[Proof of \Cref{thm:main}]
We proceed via a hybrid argument. Consider the hybrids \(G_0,G_1,\ldots,G_5\), where $G_0$ is the honest
execution and $G_5$ is Eve's attack. We say that \(G_i\) \emph{succeeds} when
its two recorded keys agree and are not \(\bot\) (an abort counts as failure).

\begin{hybrid}{Hybrid \(G_0\) (real execution)}
  \item Sample \(H_0\leftarrow\cH\).
  \item Run
  \((T_{\cA,0},M_{1,0},S_{\cA,0})
  \leftarrow\cA_1^{H_0}(1^\lambda)\).
  \item Run
  \((\Psi_0,K_{\cB,0})
  \leftarrow\cB^{H_0}(1^\lambda,M_{1,0})\).
  \item Run
  \(K_{\cA,0}
  \leftarrow\cA_2^{H_0}(1^\lambda,S_{\cA,0},\Psi_0)\).
\end{hybrid}
Therefore,
\[
  \Pr[G_0\text{ succeeds}]=\alpha.
\]

\begin{hybrid}{Hybrid \(G_1\)}
  \item Sample \(H_1\leftarrow\cH\).
  \item Run
  \((T_{\cA,1},M_{1,1},S_{\cA,1})
  \leftarrow\cA_1^{H_1}(1^\lambda)\).
  \item \changed{Run
  \(h_1\leftarrow\Learn_\delta^{H_1}(1^\lambda,M_{1,1})\); abort if
  \(h_1=\bot\).}
  \item \changed{Sample a fresh oracle
  \(\widehat H_1\leftarrow\Unif(\cH(h_1))\) and set
  \(\widetilde H_1=\widehat H_1[T_{\cA,1}]\).}
  \item Run
  \((\Psi_1,K_{\cB,1})
  \leftarrow\cB^{\changed{\widetilde H_1}}(1^\lambda,M_{1,1})\).
  \item Run
  \(K^*_1
  \leftarrow\cA_2^{\changed{\widetilde H_1}}
  (1^\lambda,S_{\cA,1},\Psi_1)\). 
\end{hybrid}

\begin{lemma}[\(G_0\) to \(G_1\)]\label{lem:g0-g1}
\[
  \left|
  \Pr[G_0\text{ succeeds}]
  -
  \Pr[G_1\text{ succeeds}]
  \right|
  \le\delta.
\]
\end{lemma}

\begin{proof}
We begin by adding an intermediate hybrid \(G'_0\) that inserts into \(G_0\) an additional call
\(h\leftarrow\Learn_\delta^{H_0}(1^\lambda,M_{1,0})\) whose output is otherwise
ignored. This computation does
not change the real execution. By
\Cref{lem:learn-success}, the call aborts with probability at most
\(\delta\).

Condition on the attacker not aborting, and fix values
\((m_1,\tau,h)\) with positive probability. By the deterministic-output
property of \(\Learn\),
\[
  \{\Learn_\delta^H(1^\lambda,m_1)=h\}
  =
  \{h\preceq H\}.
\]
Because \(\cA_1\) makes only classical oracle queries, once
\((M_1,T_{\cA})=(m_1,\tau)\) is fixed, its conditional residual state
depends only on the recorded query-answer history, not on any value of
\(H\) outside \(\Dom(\tau)\). Since the original oracle is uniform,
conditioning on
\[
  M_1=m_1,\qquad T_{\cA}=\tau,\qquad h\preceq H
\]
leaves \(H\) uniform over \(\cH(h,\tau)\) and independent of
\(S_{\cA}\). This is the classical-quantum case of the
conditional-oracle identity used by Katz and
Sela~\cite[Lemma~6]{KatzSela24}.

In \(G_1\),
\(\widehat H_1\leftarrow\Unif(\cH(h))\) is independent of Alice's conditioned
residual state, and \(\widehat H_1[\tau]\) is uniform over \(\cH(h,\tau)\). The
conditioned experiments are therefore identical. The only loss comes from
the abort event, which has probability at most \(\delta\).
\end{proof}

\begin{hybrid}{Hybrid \(G_2\)}
  \item Sample \(H_2\leftarrow\cH\).
  \item Run
  \((T_{\cA,2},M_{1,2},S_{\cA,2})
  \leftarrow\cA_1^{H_2}(1^\lambda)\).
  \item Run
  \(h_2\leftarrow\Learn_\delta^{H_2}(1^\lambda,M_{1,2})\); abort if
  \(h_2=\bot\).
  \item Sample \(\widehat H_2\leftarrow\Unif(\cH(h_2))\).
  \item \changed{Prepare
  \((T'_{\cA,2},S'_{\cA,2})\leftarrow\omega_{M_{1,2},h_2}\) and set
  \(\widetilde H_2=\widehat H_2[T'_{\cA,2}]\).}
  \item Run
  \((\Psi_2,K_{\cB,2})
  \leftarrow\cB^{\widetilde H_2}(1^\lambda,M_{1,2})\).
  \item Run
  \(K^*_2
  \leftarrow\cA_2^{\widetilde H_2}
  (1^\lambda,S'_{\cA,2},\Psi_2)\).
\end{hybrid}

\begin{lemma}[\(G_1\) to \(G_2\)]\label{lem:g1-g2}
\[
  \Pr[G_1\text{ succeeds}]
  =
  \Pr[G_2\text{ succeeds}].
\]
\end{lemma}

\begin{proof}
Both hybrids abort on the same event. Condition on a reachable
nonabort pair \((m_1,h)\). By the deterministic-output property, the
conditioned state of Alice's actual view in \(G_1\) is exactly
\(\omega_{m_1,h}\); \(G_2\) simply prepares a new instance of that same
classical-quantum state. In both hybrids, the fresh oracle
\(\widehat H_i\leftarrow\Unif(\cH(h))\) is independent of the view, and is then
reprogrammed according to its transcript. The rest of the experiments is
identical and therefore their conditional success probabilities agree for every
reachable nonabort \((m_1,h)\). Since the two hybrids have the same
distribution on \((M_1,h)\) and both fail on abort, their success probabilities
are equal.
\end{proof}

\begin{hybrid}{Hybrid \(G_3\)}
  \item Sample \(H_3\leftarrow\cH\).
  \item Run
  \((T_{\cA,3},M_{1,3},S_{\cA,3})
  \leftarrow\cA_1^{H_3}(1^\lambda)\).
  \item Run
  \(h_3\leftarrow\Learn_\delta^{H_3}(1^\lambda,M_{1,3})\); abort if
  \(h_3=\bot\).
  \item Sample \(\widehat H_3\leftarrow\Unif(\cH(h_3))\).
  \item Prepare
  \((T'_{\cA,3},S'_{\cA,3})\leftarrow\omega_{M_{1,3},h_3}\) and set
  \(\widetilde H_3=\widehat H_3[T'_{\cA,3}]\).
  \item \changed{Run
  \((\Psi_3,K_{\cB,3})
  \leftarrow\cB^{\widehat H_3}(1^\lambda,M_{1,3})\).}
  \item Run
  \(K^*_3
  \leftarrow\cA_2^{\widetilde H_3}
  (1^\lambda,S'_{\cA,3},\Psi_3)\).
\end{hybrid}

\begin{lemma}[\(G_2\) to \(G_3\)]\label{lem:g2-g3}
\[
  \left|
  \Pr[G_2\text{ succeeds}]
  -
  \Pr[G_3\text{ succeeds}]
  \right|
  \le 2q_{\cB}\sqrt\delta.
\]
\end{lemma}

\begin{proof}
Fix a reachable nonabort pair \((m_1,h)\). We construct a distinguisher
\(\cD_{m_1,h}\) for the game in \Cref{lem:reprogramming} as follows.
\begin{enumerate}[leftmargin=2em]
  \item The distinguisher samples
  \(\widehat H\leftarrow\Unif(\cH(h))\), sends
  \(F_0:=\widehat H\) to the challenger, and supplies the randomized
  algorithm \(C_{m_1,h}\). On fresh coins, \(C_{m_1,h}\) samples
  \(\tau'\leftarrow\mu_{m_1,h}\) and returns
  \[
    R:=\tau'|_{\Dom(\tau')\setminus\Dom(h)}.
  \]
  \item The challenger samples the coins of \(C_{m_1,h}\), obtains \(R\),
  sets \(F_1:=F_0[R]\), samples \(b\leftarrow\bits\), and gives the
  distinguisher quantum oracle access to \(F_b\). The distinguisher runs
  \[
    (\Psi,K_{\cB})
    \leftarrow\cB^{F_b}(1^\lambda,m_1)
  \]
  and retains \((\Psi,K_{\cB})\).
  \item The challenger removes access to \(F_b\) and reveals the coins of
  \(C_{m_1,h}\), thereby revealing \(\tau'\). The distinguisher prepares
  \(S'_{\cA}\) in the conditional state \(\sigma_{m_1,\tau'}\), runs
  \[
    K^*\leftarrow
    \cA_2^{\widehat H[\tau']}
    (1^\lambda,S'_{\cA},\Psi),
  \]
  and outputs \(1\) iff \(K^*=K_{\cB}\neq\bot\).
\end{enumerate}
Only Bob's algorithm receives the challenge oracle, so the distinguisher
makes at most \(q_{\cB}\) challenge queries.

Every transcript in \(\Supp(\mu_{m_1,h})\) is consistent with \(h\),
so \(F_1=\widehat H[R]=\widehat H[\tau']\). The cases \(b=1\) and \(b=0\)
reproduce \(G_2\) and \(G_3\), respectively, conditioned on \((m_1,h)\).
Moreover, for every \(x\),
\[
  \Pr[x\in\Dom(R)]
  \le
  \begin{cases}
    0,&x\in\Dom(h),\\
    \Pr_{\tau'\leftarrow\mu_{m_1,h}}
      [x\in\Dom(\tau')]\le\delta,&x\notin\Dom(h),
  \end{cases}
\]
where the last inequality is \Cref{lem:learn-success}. The reprogramming
lemma bounds the conditional difference by \(2q_{\cB}\sqrt\delta\).
To remove the conditioning, let \(\pi(m_1,h)\) be the common subprobability
mass of nonabort executions of the shared prefix that output
\((M_1,h)=(m_1,h)\), and write
\[
  s_i(m_1,h)
  :=
  \Pr[G_i\text{ succeeds}\mid M_{1,i}=m_1,\ h_i=h]
  \qquad (i\in\{2,3\}).
\]
The preceding argument gives
\(\lvert s_2(m_1,h)-s_3(m_1,h)\rvert
\le 2q_{\cB}\sqrt\delta\) for every \((m_1,h)\in\Supp(\pi)\).
Because success is defined to be false on abort, the two success
probabilities are the corresponding \(\pi\)-weighted sums. Hence the triangle
inequality gives
\[
  \begin{aligned}
  \left|\Pr[G_2\text{ succeeds}]-\Pr[G_3\text{ succeeds}]\right|
  &\le
  \sum_{(m_1,h)\in\Supp(\pi)}\pi(m_1,h)
    \left|s_2(m_1,h)-s_3(m_1,h)\right|\\
  &\le 2q_{\cB}\sqrt\delta,
  \end{aligned}
\]
where \(\sum_{(m_1,h)\in\Supp(\pi)}\pi(m_1,h)\le1\).
\end{proof}

\begin{hybrid}{Hybrid \(G_4\)}
  \item Sample \(H_4\leftarrow\cH\).
  \item Run
  \((T_{\cA,4},M_{1,4},S_{\cA,4})
  \leftarrow\cA_1^{H_4}(1^\lambda)\).
  \item Run
  \(h_4\leftarrow\Learn_\delta^{H_4}(1^\lambda,M_{1,4})\); abort if
  \(h_4=\bot\).
  \item \changed{Sample
  \(\widehat H_4\leftarrow\Unif(\cH(h_4))\) and set Bob's oracle
  \(H^{\mathsf B}_4=\widehat H_4[T_{\cA,4}]\).}
  \item Prepare
  \((T'_{\cA,4},S'_{\cA,4})\leftarrow\omega_{M_{1,4},h_4}\) and set
  \(H^{\mathsf E}_4=H^{\mathsf B}_4[T'_{\cA,4}]\).
  \item Run
  \((\Psi_4,K_{\cB,4})
  \leftarrow\cB^{H^{\mathsf B}_4}(1^\lambda,M_{1,4})\).
  \item Run
  \(K^*_4
  \leftarrow\cA_2^{H^{\mathsf E}_4}
  (1^\lambda,S'_{\cA,4},\Psi_4)\).
\end{hybrid}

\begin{lemma}[\(G_3\) to \(G_4\)]\label{lem:g3-g4}
\[
  \left|
  \Pr[G_3\text{ succeeds}]
  -
  \Pr[G_4\text{ succeeds}]
  \right|
  \le 2(q_{\cB}+q_{\cA,2})\sqrt\delta.
\]
\end{lemma}

\begin{proof}
Fix a reachable nonabort pair \((m_1,h)\). In \(G_4\), the original
residual register \(S_{\cA,4}\) is unused and the conditional marginal of
\(T_{\cA,4}\) is \(\mu_{m_1,h}\). Thus the conditioned hybrid is unchanged
if we independently sample
\[
  \widehat H\leftarrow\Unif(\cH(h)),
  \qquad
  \tau''\leftarrow\mu_{m_1,h},
  \qquad
  (T'_{\cA},S'_{\cA})\leftarrow\omega_{m_1,h},
\]
and then run Bob with \(\widehat H[\tau'']\) and Alice's final algorithm with
\(\widehat H[\tau''][T'_{\cA}]\). As usual, the second reprogramming
overwrites the first wherever their domains overlap.

We construct a distinguisher \(\cD_{m_1,h}\) for the game in
\Cref{lem:reprogramming}.
\begin{enumerate}[leftmargin=2em]
  \item The distinguisher prepares
  \((T'_{\cA},S'_{\cA})\leftarrow\omega_{m_1,h}\), samples
  \(\widehat H\leftarrow\Unif(\cH(h))\), and sends
  \(F_0:=\widehat H\) to the challenger. It also supplies a randomized
  algorithm \(C_{m_1,h}\) that samples
  \(\tau''\leftarrow\mu_{m_1,h}\) and returns
  \[
    R:=\tau''|_{\Dom(\tau'')\setminus\Dom(h)}.
  \]
  \item The challenger samples the coins of \(C_{m_1,h}\), obtains \(R\),
  sets \(F_1:=F_0[R]\), samples \(b\leftarrow\bits\), and gives the
  distinguisher quantum oracle access to \(F_b\). The distinguisher runs
  \[
    (\Psi,K_{\cB})\leftarrow\cB^{F_b}(1^\lambda,m_1)
  \]
  followed by
  \[
    K^*\leftarrow
    \cA_2^{F_b[T'_{\cA}]}
    (1^\lambda,S'_{\cA},\Psi),
  \]
  and records the bit \(d:=\mathbf 1[K^*=K_{\cB}\neq\bot]\).
  \item The challenger removes access to \(F_b\) and reveals the coins of
  \(C_{m_1,h}\), thereby revealing \(\tau''\). The distinguisher outputs
  the recorded bit \(d\).
\end{enumerate}

Each query to \(F_b[T'_{\cA}]\) uses one query to \(F_b\). First compute
reversibly whether the query point lies in \(\Dom(T'_{\cA})\). On that
branch, swap the answer register with an ancilla in
\(\ket{+}^{\otimes\ell}\), apply the XOR oracle \(U_{F_b}\), and swap back.
Because every XOR translation fixes \(\ket{+}^{\otimes\ell}\), the challenge
oracle has no effect on the reprogrammed branch.  XOR the known value
\(T'_{\cA}(x)\) into the answer register there and uncompute the lookup.  This
implements \(U_{F_b[T'_{\cA}]}\) exactly. The total number of challenge queries
is at most \(q_{\cB}+q_{\cA,2}\).

If \(b=0\), this is \(G_3\). If \(b=1\), consistency on
\(\Dom(h)\) gives \(F_1=\widehat H[R]=\widehat H[\tau'']\), yielding the
equivalent form of \(G_4\). By \Cref{lem:learn-success},
\(\max_x\Pr[x\in\Dom(R)]\le\delta\). The reprogramming lemma gives the
pointwise bound
\[
  \left|
    \Pr[G_3\text{ succeeds}\mid M_1=m_1,\ h_3=h]
    -
    \Pr[G_4\text{ succeeds}\mid M_1=m_1,\ h_4=h]
  \right|
  \le 2(q_{\cB}+q_{\cA,2})\sqrt\delta
\]
for every reachable nonabort pair \((m_1,h)\). Let \(\pi(m_1,h)\) denote the
common subprobability mass of nonabort executions of the shared prefix that
produce that pair, and define
\[
  s_i(m_1,h)
  :=
  \Pr[G_i\text{ succeeds}\mid M_{1,i}=m_1,\ h_i=h]
  \qquad (i\in\{3,4\}).
\]
Because success is defined to be false on abort, the two success
probabilities are the corresponding \(\pi\)-weighted sums. The triangle
inequality yields
\[
  \begin{aligned}
  \left|\Pr[G_3\text{ succeeds}]-\Pr[G_4\text{ succeeds}]\right|
  &\le
  \sum_{(m_1,h)\in\Supp(\pi)}\pi(m_1,h)
    \left|s_3(m_1,h)-s_4(m_1,h)\right|\\
  &\le 2(q_{\cB}+q_{\cA,2})\sqrt\delta.
  \end{aligned}
\]
\end{proof}

\begin{hybrid}{Hybrid \(G_5\)}
  \item Sample \(H_5\leftarrow\cH\).
  \item Run
  \((T_{\cA,5},M_{1,5},S_{\cA,5})
  \leftarrow\cA_1^{H_5}(1^\lambda)\).
  \item Run
  \(h_5\leftarrow\Learn_\delta^{H_5}(1^\lambda,M_{1,5})\); abort if
  \(h_5=\bot\).
  \item Prepare
  \((T'_{\cA,5},S'_{\cA,5})\leftarrow\omega_{M_{1,5},h_5}\) and set
  \(\widetilde H_5=H_5[T'_{\cA,5}]\).
  \item Run
  \((\Psi_5,K_{\cB,5})
  \leftarrow\cB^{H_5}(1^\lambda,M_{1,5})\).
  \item Run
  \(K^*_5
  \leftarrow\cA_2^{\widetilde H_5}
  (1^\lambda,S'_{\cA,5},\Psi_5)\).
\end{hybrid}

\begin{lemma}[\(G_4\) to \(G_5\)]\label{lem:g4-g5}
\[
  \Pr[G_4\text{ succeeds}]
  =
  \Pr[G_5\text{ succeeds}].
\]
\end{lemma}

\begin{proof}
Both hybrids abort on the same event. On success, fix
\((m_1,\tau,h)\). By the deterministic-output property and the
conditional-oracle identity of Katz and Sela~\cite[Lemma~6]{KatzSela24},
conditioned on this triple the
original oracle \(H_5\) in \(G_5\) is uniform over \(\cH(h,\tau)\) and
independent of the unused original residual state. In \(G_4\), sampling
\(\widehat H_4\leftarrow\Unif(\cH(h))\) and setting
\(H^{\mathsf B}_4=\widehat H_4[\tau]\) gives the
same uniform distribution over \(\cH(h,\tau)\). The independently prepared
view \(\omega_{m_1,h}\) and every subsequent operation are the same.
Thus the conditional success probabilities of \(G_4\) and \(G_5\) are equal
for every reachable nonabort triple \((m_1,\tau,h)\). The two hybrids induce
the same distribution \(\pi(m_1,\tau,h)\) on these triples. Define
\[
  s_i(m_1,\tau,h)
  :=
  \Pr[G_i\text{ succeeds}\mid
    M_{1,i}=m_1,\ T_{\cA,i}=\tau,\ h_i=h]
  \qquad (i\in\{4,5\}).
\]
Then \(s_4(m_1,\tau,h)=s_5(m_1,\tau,h)\) throughout \(\Supp(\pi)\).
Both hybrids have success probability zero on abort, so the law of total
probability gives
\[
  \Pr[G_4\text{ succeeds}]
  =
  \sum_{m_1,\tau,h}\pi(m_1,\tau,h)
    s_4(m_1,\tau,h)
  =
  \sum_{m_1,\tau,h}\pi(m_1,\tau,h)
    s_5(m_1,\tau,h)
  =
  \Pr[G_5\text{ succeeds}].
\]
\end{proof}

Hybrid \(G_5\) is exactly \(\IntKR_{\Pi,\cE_\delta}\): the preprocessing
stage learns \(h\) and prepares the conditioned view; Bob runs with the
original oracle; Eve intercepts \(\Psi_5\); and Eve runs Alice's final
algorithm with \(H_5[T'_{\cA,5}]\).  Therefore
\begin{align*}
  \Adv^{\IntKR}_{\Pi,\cE_\delta}(\lambda)
  &=
  \Pr[G_5\text{ succeeds}]\\
  &\ge
  \Pr[G_0\text{ succeeds}]
  -\delta
  -2q_{\cB}\sqrt\delta
  -2(q_{\cB}+q_{\cA,2})\sqrt\delta\\
  &=
  \alpha-\delta-2(2q_{\cB}+q_{\cA,2})\sqrt\delta.
\end{align*}
The query bound follows from \Cref{def:eve}.
\end{proof}

\begin{corollary}[Concrete parameter choice]\label{cor:parameters}
Let \(\Pi\) satisfy \Cref{def:protocol}, let
\(\alpha=\alpha_\Pi(\lambda)>0\), and define
\[
  Q:=\max\{1,2q_{\cB}+q_{\cA,2}\},
  \qquad
  \delta:=\left(\frac{\alpha}{8Q}\right)^2.
\]
Then the adversary of \Cref{thm:main} succeeds with probability strictly
greater than \(\alpha/2\) and makes at most
\[
  \ceil{4096\,q_{\cA,1}Q^4/\alpha^4}+q_{\cA,2}
\]
random-oracle queries.
\end{corollary}

\begin{proof}
The reprogramming loss is
\[
  2Q\sqrt\delta=\alpha/4.
\]
Since \(0<\alpha\le1\) and \(Q\ge1\),
\(\delta=\alpha^2/(64Q^2)\le\alpha/64\).  Hence the success lower bound is
at least
\[
  \alpha-\alpha/64-\alpha/4>\alpha/2.
\]
Moreover,
\[
  \frac{q_{\cA,1}}{\delta^2}
  =
  \frac{4096\,q_{\cA,1}Q^4}{\alpha^4}.
\]
\end{proof}

\begin{corollary}[Polynomial-query two-message impossibility]\label{cor:quantum-impossibility}
Let \(\Pi=\{\Pi_\lambda\}\) be a family of protocols satisfying
\Cref{def:protocol}. Suppose that \(q_{\cA,1}\), \(q_{\cB}\), and
\(q_{\cA,2}\) are each bounded by \(\poly(\lambda)\), and that
\(\alpha_\Pi(\lambda)\ge1/\poly(\lambda)\). Then the family is not
interception key-recovery secure against computationally unbounded quantum
adversaries making at most \(\poly(\lambda)\) random-oracle queries. In
particular, the conclusion applies when the completeness error is negligible.
\end{corollary}

\subsection{The Multi-Round Attack}
\label{sec:multi-round}

We give an attack on a multi-round key agreement protocol with a classical prefix and one final possibly quantum message, as defined in 
\Cref{def:multi-protocol}. Recall that in this model, the only property the prefix must satisfy is that all communication and all queries to the random oracle are classical;
the parties' local computation and residual states may otherwise be
quantum.  Our attack and corresponding analysis extend the techniques from \cite{BarakMahmoodyFull} to our setting.

\subsubsection{Classical Views of the Prefix}
\label{sec:classical-prefix-views}

The main difference between our setting and the purely classical setting in \cite{BarakMahmoodyFull} is that we allow the parties' local computation to be quantum. However, as we discuss in this section, we can assume without loss of generality that Alice and Bob are executing classical algorithms to generate the prefix. In turn, we will be able to apply the techniques from \cite{BarakMahmoodyFull} directly. 

We first explain how to generate classical views of the prefix. Let $\Pi$ be a protocol as defined in Definition~\ref{def:multi-protocol}. Recall that while the parties must be query bounded, they may be computationally unbounded. We therefore can define an exact classical
simulation of the prefix of $\Pi$ as follows. The simulator maintains a classical descirption of the parties' residual quantum states. Whenever an honest party sends a
message or queries the oracle during the prefix, the simulator produces a sample according to the distribution induced by the party's current
quantum state and updates its classical description of the corresponding
conditional residual state. Thus the simulation makes the same classical
oracle queries and produces exactly the same joint state of the public
transcript, the parties' query transcripts, and their residual registers as
the original prefix.

For $P\in\{\cA,\cB\}$, define
\[
  U_P := (R_P,M,T_P)
\]
as the classical view of party $P$ in this simulation, where $R_P$ records all
private random choices of the simulator, $M$ is the public prefix transcript,
and $T_P$ is the complete ordered oracle query-answer transcript of $P$.
This is equivalent to the notion of a view used by Barak and
Mahmoody~\cite[Section~3.1]{BarakMahmoodyFull}.  We note that the classical description of the residual states is not part of $U_P$.

We now explain what is needed to use this simulated classical view in combination with the techniques of \cite{BarakMahmoodyFull}. The eventual attack will retain the public prefix transcript $m$ and a
partial oracle $h$ learned during the prefix, and then prepare a fresh Alice
prefix view from its conditional marginal given $(m,h)$.  To justify replacing
Alice's real view by this fresh sample without significantly changing Bob's
subsequent behavior, we need the real joint state of the two prefix views,
conditioned on $(m,h)$, to be close to the product of its marginals.
\Cref{sec:bm-learner} obtains the corresponding approximate-product
guarantee for the simulated classical views $U_{\cA}$ and $U_{\cB}$ from
the Barak--Mahmoody learner.

That classical guarantee is not yet sufficient for our application.  The
final round of the key agreement protocol acts on the parties' residual quantum registers, rather than on the
simulator's private randomness, and its algorithms may query the same random
oracle $H$ in superposition.  We must therefore both transfer the
approximate-product guarantee from $U_{\cA},U_{\cB}$ to the parties' actual
query transcripts and residual quantum states, and preserve the correct
correlation of those states with $H$.  The local channels defined below
perform the first task.  The conditional-oracle identity in
\Cref{eq:prefix-conditional-oracle-main} performs the second.  In
\Cref{sec:conditioned-prefix-resampling}, we combine the two facts to replace
the real conditioned prefix state by independently sampled Alice and Bob
views together with a suitably sampled and reprogrammed oracle.

For every $u_P$ in the support of $U_P$, let $T_P(u_P)$ denote the query
transcript contained in $u_P$, and let $\sigma_P(u_P)$ be the residual density
operator tracked by the simulator after the prefix on that view.  Define the
local channel
\[
  \Phi_P(\proj{u_P})
  :=
  \proj{T_P(u_P)}\otimes\sigma_P(u_P).
\]
Thus $U_P$ is only an auxiliary classical view used to invoke the
Barak--Mahmoody result; applying $\Phi_P$ discards the simulator's private
randomness and keeps only the actual query transcript and residual quantum
state.

For a complete public prefix transcript $m$ and a partial oracle $h$, write
\[
  \mathsf F_{m,h}:=\{M=m,\ h\preceq H\}.
\]
Whenever $\Pr[\mathsf F_{m,h}]>0$, let $\nu^{AB}_{m,h}$ be the diagonal density
operator corresponding to the conditional distribution of
$(U_{\cA},U_{\cB})$ given $\mathsf F_{m,h}$, and let
$\nu^A_{m,h},\nu^B_{m,h}$ be its marginals.  Define
\[
  \omega^{AB}_{m,h}
  :=
  (\Phi_{\cA}\otimes\Phi_{\cB})(\nu^{AB}_{m,h}),
\]
and define $\omega^A_{m,h},\omega^B_{m,h}$ analogously.  Hence
$\omega^{AB}_{m,h}$ is a state on
$T_{\cA}S_{\cA}T_{\cB}S_{\cB}$ and contains no simulator randomness.
We note that, since we condition on $M=m$, the transcript register is fixed to
$\proj{m}_M$.  Let
$\Omega_{m,h}$ be the real
conditional state on
$\mathsf H T_{\cA}S_{\cA}T_{\cB}S_{\cB}$ given $\mathsf F_{m,h}$, where
$\mathsf H$ stores the random oracle.  Since the simulation is exact,
$\omega^{AB}_{m,h}$ is precisely the
$T_{\cA}S_{\cA}T_{\cB}S_{\cB}$-marginal of $\Omega_{m,h}$.

Since all oracle queries in the prefix are classical, once
$M=m$, $T_{\cA}=t_{\cA}$, and $T_{\cB}=t_{\cB}$ are fixed, the
residual state of Alice and Bob is the same for every oracle consistent
with $t_{\cA}$ and $t_{\cB}$.  Therefore, since $H$ is initially
uniform, whenever the conditioning event has positive probability,
\begin{equation}
  H\ \big|\
  \bigl(
    \mathsf F_{m,h},
    T_{\cA}=t_{\cA},
    T_{\cB}=t_{\cB}
  \bigr)
  \ \sim\
  \Unif\!\left(
    \cH(h,t_{\cA},t_{\cB})
  \right).
  \label{eq:prefix-conditional-oracle-main}
\end{equation}
Moreover, under this conditioning, the residual state is independent
of $H$.

\subsubsection{The Barak--Mahmoody Learner}
\label{sec:bm-learner}

In this section, we recall the relevant results from \cite{BarakMahmoodyFull}, lifted to our setting.  The first lemma below restates the
general result of Barak and Mahmoody for an arbitrary classical protocol.  In
particular, their Theorem~4.3 gives an expected query bound and guarantees
approximate independence and lightness after each round.  We retain only the
guarantees after the last protocol message, when \(M\) is the complete public
transcript and \(h\) is the learner's final partial oracle.  The second lemma
then applies this result to the exact classical simulation from
\Cref{sec:classical-prefix-views}.  It fixes the parameters in terms of
\(\eta\), truncates the learner to obtain a worst-case query bound, aborts
if the truncation threshold is reached or either the independence or
lightness guarantee fails, and transfers the
approximate-independence guarantee to the parties' residual quantum states.
It also characterizes the learner's output by \(h\preceq H\).  Thus the first
lemma isolates the result imported from Barak and Mahmoody, while the second
puts it in the form used for conditioned prefix resampling in
\Cref{sec:conditioned-prefix-resampling}.  The Barak--Mahmoody statement uses
Items~2, 4, and~5 of their Theorem~4.3 and their definition of
self-dependency~\cite[Definition~4.2 and Theorem~4.3,
Items~2, 4, and~5]{BarakMahmoodyFull}.

\begin{lemma}[Barak--Mahmoody]
\label{lem:bm-final-transcript}
Consider a classical two-party random-oracle protocol.  Let $M$ be its complete
public transcript and let
\[
  U_P=(R_P,M,T_P),\qquad P\in\{\cA,\cB\},
\]
be the parties' classical views, where $T_P$ is the oracle query-answer
transcript.  Suppose Alice and Bob make at most $n_A$ and $n_B$ oracle
queries, respectively, where $n_A,n_B\ge1$.

Let
\[
  0<\varepsilon_1\le\varepsilon_2<1/10,
  \qquad
  0<\alpha,\beta<1,
  \qquad
  \alpha\beta\ge\varepsilon_2.
\]
There is a deterministic classical-query algorithm
\[
  \mathsf L^H_{\varepsilon_1,\varepsilon_2}(m)
\]
that takes a complete public transcript $m$ as input and outputs a partial
oracle $h$.  For every realized pair $(m,h)$, let $\mu^{AB}_{m,h}$ be the
conditional distribution of $(U_{\cA},U_{\cB})$ given
\[
  M=m,
  \qquad
  \mathsf L^H_{\varepsilon_1,\varepsilon_2}(m)=h,
\]
and let $\mu^A_{m,h},\mu^B_{m,h}$ be its marginals.  Then:
\begin{enumerate}[leftmargin=2em,label=\textup{(\roman*)}]
  \item The expected number of oracle queries made by
  $\mathsf L_{\varepsilon_1,\varepsilon_2}$ is at most
  \[
    \frac{4n_A n_B}{\varepsilon_1}.
  \]

  \item Except with probability at most $9\alpha$ over the realized pair
  $(M,h)$,
  \[
    \Delta_{\mathrm{TV}}\!\left(
      \mu^{AB}_{M,h},
      \mu^A_{M,h}\times\mu^B_{M,h}
    \right)
    \le 9\beta.
  \]

  \item Except with probability at most $9\alpha$ over the realized pair
  $(M,h)$, every $x\notin\Dom(h)$ satisfies
  \[
    \Pr[x\in\Dom(T_{\cA})\mid M,h]
      < \frac{\varepsilon_2}{n_B}+\beta,
    \qquad
    \Pr[x\in\Dom(T_{\cB})\mid M,h]
      < \frac{\varepsilon_2}{n_A}+\beta.
  \]
\end{enumerate}
\end{lemma}

\begin{proof}
We use the deterministic learner from Barak and Mahmoody's
Construction~3.3 with parameter $\varepsilon_1$, together with the stopping
modification from their Theorem~4.3.  Barak and Mahmoody describe this learner
as running alongside the protocol.  Its oracle queries do not affect the
honest execution, and Construction~3.3 uses deterministic tie-breaking.
Consequently, given the complete transcript $m$, we may replay the learner on
the successive prefixes of $m$ and make the same oracle queries.  This defines
the deterministic algorithm $\mathsf L^H_{\varepsilon_1,\varepsilon_2}(m)$
with the same final partial oracle and query complexity.

With the stopping modification, the learner is an
$(\varepsilon_1,\varepsilon_2)$-attacker in the sense of Definition~4.1 of
Barak and Mahmoody.  Item~2 of their Theorem~4.3 gives the expected query
bound in Item~\textup{(i)}.  Their Definition~4.2 identifies self-dependency
with statistical distance from the product of the marginals.  Evaluating
Item~4 of Theorem~4.3 after the final protocol message therefore gives
Item~\textup{(ii)}.  Evaluating Item~5 at the same point gives the two
lightness bounds in Item~\textup{(iii)}
\cite[Construction~3.3, Definitions~4.1--4.2, and Theorem~4.3,
Items~2, 4, and~5]{BarakMahmoodyFull}.
\end{proof}

\begin{lemma}[Barak--Mahmoody prefix attacker]
\label{lem:multi-independence-attacker}
Let $0<\eta<1/100$, and set
\[
  \bar q_{\cA}:=\max\{1,q_{\cA,\mathsf{pre}}\},
  \qquad
  \bar q_{\cB}:=\max\{1,q_{\cB,\mathsf{pre}}\},
\]
\[
  Q_{\mathsf{learn}}
  :=
  \left\lceil
    \frac{4\cdot 10^5\,\bar q_{\cA}\bar q_{\cB}}{\eta^3}
  \right\rceil.
\]
There is a deterministic algorithm
$\Learn^{\mathsf{BM},H}_\eta$ with classical oracle access that, on input a
complete public prefix transcript $m$, outputs a partial oracle $h$ or aborts.
A pair $(m,h)$ is \emph{reachable} if
\[
  \Pr\!\left[
    M=m,\
    \Learn^{\mathsf{BM},H}_\eta(1^\lambda,m)=h
  \right]>0.
\]
It satisfies the following properties.
\begin{enumerate}[leftmargin=2em,label=\textup{(\roman*)}]
  \item It makes at most $Q_{\mathsf{learn}}$ oracle queries and aborts with
  probability at most $\eta$.

  \item For every reachable pair $(m,h)$,
  \[
    \frac12
    \left\|
      \omega^{AB}_{m,h}
      -
      \omega^A_{m,h}\otimes\omega^B_{m,h}
    \right\|_1
    \le\eta.
  \]

  \item For every reachable pair $(m,h)$ and every $x\notin\Dom(h)$,
  \[
    \Pr[x\in\Dom(T_{\cA})\mid\mathsf F_{m,h}]
    \le\eta.
  \]
\end{enumerate}
Moreover, for every reachable pair $(m,h)$ and every oracle $H$,
\[
  \Learn^{\mathsf{BM},H}_\eta(1^\lambda,m)=h
  \quad\Longleftrightarrow\quad
  h\preceq H.
\]
\end{lemma}

\begin{proof}
We apply \Cref{lem:bm-final-transcript} to the exact classical simulation of
the prefix with parameters
\[
  n_A=\bar q_{\cA},
  \qquad
  n_B=\bar q_{\cB},
  \qquad
  \varepsilon_1=\varepsilon_2=\frac{\eta^2}{10^4},
  \qquad
  \alpha_0=\beta_0=\frac{\eta}{100}.
\]
These parameters satisfy $\alpha_0\beta_0=\varepsilon_2$.  The resulting
learner makes at most
\[
  \frac{4\bar q_{\cA}\bar q_{\cB}}{\varepsilon_1}
  =
  \frac{4\cdot10^4\bar q_{\cA}\bar q_{\cB}}{\eta^2}
\]
queries in expectation.  We truncate this learner after
$Q_{\mathsf{learn}}$ queries and abort if the truncation threshold is
reached.  By Markov's inequality, this occurs with probability at most
$\eta/10$.

The exceptional events in Items~\textup{(ii)} and~\textup{(iii)} of
\Cref{lem:bm-final-transcript} have total probability at most
$18\alpha_0$.  Since computation is unbounded in our query model, once the
learner terminates we may test whether both conditional guarantees hold and
abort otherwise, without making any further oracle queries.  By a union
bound, the total abort probability is at most
\[
  18\alpha_0+\frac{\eta}{10}
  =\frac{28\eta}{100}<\eta.
\]
This proves Item~\textup{(i)}.

We next establish the claimed characterization of the learner's output.  Fix
a reachable nonabort pair $(m,h)$.  For fixed $m$, every query and stopping
decision of the learner is a deterministic function of the partial oracle
learned so far.  Since some execution produces $h$, every oracle extending
$h$ gives the same answers along that execution and therefore produces the
same output.  Conversely, if the learner outputs $h$, then every pair in $h$
was obtained from an oracle query, and hence $h\preceq H$.  It follows that,
for every oracle $H$,
\[
  \Learn^{\mathsf{BM},H}_\eta(1^\lambda,m)=h
  \quad\Longleftrightarrow\quad
  h\preceq H.
\]
Thus, for every reachable nonabort pair $(m,h)$, conditioning on the learner's
output is equivalent to conditioning on $\mathsf F_{m,h}$.

It remains to transfer the guarantees for the simulated classical views to
the actual prefix views.  By Item~\textup{(ii)} of
\Cref{lem:bm-final-transcript},
\[
  \Delta_{\mathrm{TV}}\!\left(
    \nu^{AB}_{m,h},
    \nu^A_{m,h}\times\nu^B_{m,h}
  \right)
  \le 9\beta_0<\eta.
\]
Applying the local channel $\Phi_{\cA}\otimes\Phi_{\cB}$ and using
trace-distance data processing~\cite[Chapter~3]{Watrous18} proves
Item~\textup{(ii)}.  Finally, Item~\textup{(iii)} of
\Cref{lem:bm-final-transcript} gives, for every $x\notin\Dom(h)$,
\[
  \Pr[x\in\Dom(T_{\cA})\mid\mathsf F_{m,h}]
  <
  \frac{\varepsilon_2}{\bar q_{\cB}}+\beta_0
  \le
  \frac{\eta^2}{10^4}+\frac{\eta}{100}
  <\eta,
\]
which proves Item~\textup{(iii)}.
\end{proof}

\subsubsection{Conditioned Prefix Resampling}
\label{sec:conditioned-prefix-resampling}

The preceding subsection shows that, conditioned on a reachable pair
\((m,h)\), the parties' actual prefix views are close to the product state
\(\omega^A_{m,h}\otimes\omega^B_{m,h}\).  This does not by itself justify
sampling the two views independently in the attack, because that statement
does not include the random oracle.  Bob's final algorithm and Eve's
simulation of Alice both query the same oracle in superposition, so we must
also specify how the independently sampled views are coupled to that oracle.

The conditional-oracle identity in
\Cref{eq:prefix-conditional-oracle-main} tells us how to do so.  After
sampling the two marginal views independently, we first sample a base oracle
\(H_0\) consistent with \(h\) and Bob's transcript \(T'_{\cB}\), and then
reprogram it on Alice's transcript \(T'_{\cA}\).  The order is deliberate:
the base oracle is compatible with Bob's sampled view, while the final
oracle \(\widetilde H=H_0[T'_{\cA}]\) contains the answers recorded in
Alice's sampled view. 

The proof compares the real and resampled experiments by applying the same
channel to two input states.  On the real joint view
\(\omega^{AB}_{m,h}\), the channel reconstructs the real conditioned state
\(\Omega_{m,h}\) exactly.  On the product of the marginals, it produces the
resampled state \(\widetilde\Omega_{m,h}\) defined below.  The trace-distance
bound then essentially follows directly from the
approximate-independence guarantee of the preceding subsection.  In
particular, any later joint quantum computation using the oracle and the
prefix views changes the success probability by at most \(\eta\) when the
real conditioned prefix is replaced by this resampled one.

\begin{lemma}[Conditioned prefix resampling]
\label{lem:prefix-resampling}
Fix a reachable pair $(m,h)$, and consider the following experiment:
\begin{enumerate}[leftmargin=2em]
  \item independently prepare
  \[
    (T'_{\cA},S'_{\cA})\leftarrow\omega^A_{m,h},
    \qquad
    (T'_{\cB},S'_{\cB})\leftarrow\omega^B_{m,h};
  \]
  \item sample $H_0\leftarrow\Unif(\cH(h,T'_{\cB}))$, and define
  \[
    \widetilde H=H_0[T'_{\cA}].
  \]
\end{enumerate}
Let $\widetilde\Omega_{m,h}$ be the resulting state on
$\mathsf H T'_{\cA}S'_{\cA}T'_{\cB}S'_{\cB}$, with $\mathsf H$ holding the
sampled function $\widetilde H$.  Then
\[
  \frac12
  \left\|
    \Omega_{m,h}
    -
    \widetilde\Omega_{m,h}
  \right\|_1
  \le \eta.
\]
\end{lemma}

\begin{proof}
We compare the two states by applying the same channel to the real joint
prefix view and to the product of its marginals.  Define a channel
$\mathcal R_{m,h}$ on $T_{\cA}S_{\cA}T_{\cB}S_{\cB}$ as follows.  On
classical transcripts $t_{\cA},t_{\cB}$, it samples
\[
  H_0\leftarrow\Unif(\cH(h,t_{\cB}))
\]
and attaches the oracle $H_0[t_{\cA}]$, while leaving the four input
registers unchanged.  Every $t_{\cB}$ in the support of
$\omega^B_{m,h}$ is consistent with $h$, so this channel is well defined on
both input states considered below.

We first apply the channel to the real joint prefix view.  In this case,
$t_{\cA}$ and $t_{\cB}$ are mutually consistent, and
$H_0[t_{\cA}]$ is uniform over $\cH(h,t_{\cA},t_{\cB})$.  By
\eqref{eq:prefix-conditional-oracle-main}, this is precisely the conditional
distribution of the real oracle given the two query transcripts.  Moreover,
the residual state is independent of the unobserved oracle values under this
conditioning.  Therefore,
\[
  \mathcal R_{m,h}(\omega^{AB}_{m,h})=\Omega_{m,h}.
\]
Applying the same channel to the product of the two marginal views gives,
by the definition of the resampling experiment,
\[
  \mathcal R_{m,h}
  (\omega^A_{m,h}\otimes\omega^B_{m,h})
  =\widetilde\Omega_{m,h}.
\]
Combining these two identities with trace-distance data processing and
\Cref{lem:multi-independence-attacker}, we obtain
\[
  \frac12\left\|
    \Omega_{m,h}-\widetilde\Omega_{m,h}
  \right\|_1
  \le
  \frac12\left\|
    \omega^{AB}_{m,h}
    -\omega^A_{m,h}\otimes\omega^B_{m,h}
  \right\|_1
  \le\eta.
\]
\end{proof}

\subsubsection{The Multi-Round attacker}

In this section we conclude the proof by defining the attacker.

\begin{definition}[Multi-round reprogramming attack]
\label{def:multi-eve}
Fix \(0<\eta<1/100\). The attacker \(\cE_\eta\) acts as follows. During the
classical prefix, Eve runs the online
attacker \(\Learn^{\mathsf{BM},H}_\eta\) of
\Cref{lem:multi-independence-attacker}. If the attacker aborts, Eve records an
abort flag. Otherwise, after the last classical message Eve has the full
public transcript \(m\) and a partial oracle \(h\). She prepares an
independent Alice prefix view
\[
  V'_{\cA}=(W'_{\cA},S'_{\cA})
  \leftarrow\omega^A_{m,h},
  \qquad
  T'_{\cA}:=T_{\cA}(W'_{\cA}).
\]
After intercepting Bob's final message register \(\Psi\), Eve outputs
\(\bot\) if the attacker aborted. Otherwise she runs
\[
  K^*
  \leftarrow
  \cA_{\mathsf{fin}}^{H[T'_{\cA}]}
  (1^\lambda,m,S'_{\cA},\Psi)
\]
and outputs \(K^*\).
\end{definition}

\begin{theorem}[Multi-round quantitative attack]
\label{thm:multi-main}
Let \(\Pi^{\mathsf{mr}}\) satisfy
\Cref{def:multi-protocol}, and let
\(\alpha=\alpha_{\Pi^{\mathsf{mr}}}(\lambda)\).  For every
\(0<\eta<1/100\), the attacker of \Cref{def:multi-eve} satisfies
\[
  \Adv^{\IntKR,\mathsf{mr}}_{
    \Pi^{\mathsf{mr}},\cE_\eta}(\lambda)
  \ge
  \alpha
  -5\eta
  -2(2q_{\cB,\mathsf q}+q_{\cA,\mathsf{fin}})\sqrt\eta.
\]
It makes at most
\[
  \left\lceil
    \frac{4\cdot 10^5\,
      \max\{1,q_{\cA,\mathsf{pre}}\}
      \max\{1,q_{\cB,\mathsf{pre}}\}}
      {\eta^3}
  \right\rceil
  +q_{\cA,\mathsf{fin}}
\]
random-oracle queries.  The first term consists only of classical queries.
There is no explicit dependence on the number of prefix rounds beyond the
two parties' total prefix-query bounds.
\end{theorem}

\begin{proof}
We proceed via a hybrid argument. Hybrid \(J_0\) is the honest execution. In \(J_1,\ldots,J_5\), the online
attacker runs alongside the real classical prefix and an abort counts as
failure. Unless stated otherwise, these hybrids induce the same distribution
on the public transcript \(M\) and the attacker's output \(h\). We say that
\(J_i\) succeeds when its two recorded keys agree and are not \(\bot\).
Throughout, \(V_P=(W_P,S_P)\) denotes the extended classical-quantum
prefix view defined above. Changes from one hybrid to the next are shown in
\mrchanged{color}.

\begin{hybrid}{Multi-round hybrid \(J_0\) (real execution)}
  \item Sample \(H_0\leftarrow\cH\) and run the honest classical prefix,
  obtaining
  \((M^{(0)},V_{\cA,0},V_{\cB,0})\).
  \item Run
  \((\Psi_0,K_{\cB,0})\leftarrow
    \cB_{\mathsf q}^{H_0}(1^\lambda,M^{(0)},S_{\cB,0})\).
  \item Run
  \(K_{\cA,0}\leftarrow
    \cA_{\mathsf{fin}}^{H_0}
    (1^\lambda,M^{(0)},S_{\cA,0},\Psi_0)\).
\end{hybrid}
Thus
\nopagebreak[4]
\[
  \Pr[J_0\text{ succeeds}]=\alpha.
\]

\begin{hybrid}{Multi-round hybrid \(J_1\)}
  \item \mrchanged{Sample \(H_1\leftarrow\cH\), and run the honest
  classical prefix using \(H_1\) while activating
  \(\Learn^{\mathsf{BM},H_1}_\eta\) after every message.  Obtain
  \((M^{(1)},V_{\cA,1},V_{\cB,1},h_1)\), and abort if the attacker aborts.}
  \item Run
  \((\Psi_1,K_{\cB,1})\leftarrow
    \cB_{\mathsf q}^{H_1}(1^\lambda,M^{(1)},S_{\cB,1})\), followed by
  \(K_{\cA,1}\leftarrow
    \cA_{\mathsf{fin}}^{H_1}
    (1^\lambda,M^{(1)},S_{\cA,1},\Psi_1)\).
\end{hybrid}
Running the attacker alongside the protocol does not change the honest
execution. The only difference is that \(J_1\) counts a attacker abort as
failure. Item~(i) of
\Cref{lem:multi-independence-attacker} gives
\begin{equation}
\label{eq:multi-j0-j1}
  \left|
    \Pr[J_0\text{ succeeds}]-\Pr[J_1\text{ succeeds}]
  \right|
  \le\eta.
\end{equation}

\begin{hybrid}{Multi-round hybrid \(J_2\)}
  \item Run the real prefix and attacker as in \(J_1\), retaining
  \((M^{(2)},h_2)\) and discarding the original oracle and prefix views.
  \item \mrchanged{Independently prepare
  \[
    V'_{\cA,2}=(W'_{\cA,2},S'_{\cA,2})
      \leftarrow\omega^A_{M^{(2)},h_2},
  \]
  \[
    V'_{\cB,2}=(W'_{\cB,2},S'_{\cB,2})
      \leftarrow\omega^B_{M^{(2)},h_2}.
  \]
  Set \(T'_{\cA,2}=T_{\cA}(W'_{\cA,2})\) and
  \(T'_{\cB,2}=T_{\cB}(W'_{\cB,2})\).}
  \item \mrchanged{Sample
  \(\widehat H_2\leftarrow\Unif(\cH(h_2,T'_{\cB,2}))\) and set
  \(\widetilde H_2=\widehat H_2[T'_{\cA,2}]\).}
  \item Run
  \((\Psi_2,K_{\cB,2})\leftarrow
    \cB_{\mathsf q}^{\widetilde H_2}
    (1^\lambda,M^{(2)},S'_{\cB,2})\).
  \item Run
  \(K^*_2\leftarrow
    \cA_{\mathsf{fin}}^{\widetilde H_2}
    (1^\lambda,M^{(2)},S'_{\cA,2},\Psi_2)\).
\end{hybrid}

Condition on a reachable pair \((m,h)\). By the output characterization in
\Cref{lem:multi-independence-attacker}, the oracle and prefix views in \(J_1\)
have joint state \(\Omega_{m,h}\), while the resampled objects in \(J_2\)
have state \(\widetilde\Omega_{m,h}\). Applying
\Cref{lem:prefix-resampling} and trace-distance data processing through the
final algorithms~\cite[Chapter~3]{Watrous18} shows that, for every reachable
nonabort pair,
\[
  \left|
    \Pr[J_1\text{ succeeds}\mid M=m,\ h_1=h]
    -
    \Pr[J_2\text{ succeeds}\mid M=m,\ h_2=h]
  \right|
  \le 2\eta.
\]
Let \(\pi(m,h)\) be the common subprobability mass of nonabort prefix
executions in the two hybrids that output \((M,h)=(m,h)\). Since success is
defined to be false on abort, only these executions contribute. Define
\[
  s_i(m,h)
  :=
  \Pr[J_i\text{ succeeds}\mid M^{(i)}=m,\ h_i=h]
  \qquad (i\in\{1,2\}).
\]
The triangle inequality gives
\[
  \begin{aligned}
  &\left|
    \Pr[J_1\text{ succeeds}]-\Pr[J_2\text{ succeeds}]
  \right|\\
  &\qquad\le
  \sum_{(m,h)\in\Supp(\pi)}\pi(m,h)
  \left|s_1(m,h)-s_2(m,h)\right|
  \le 2\eta.
  \end{aligned}
\]
In particular,
\begin{equation}
\label{eq:multi-j1-j2}
  \left|
    \Pr[J_1\text{ succeeds}]-\Pr[J_2\text{ succeeds}]
  \right|
  \le2\eta.
\end{equation}

\begin{hybrid}{Multi-round hybrid \(J_3\)}
  \item Run the real prefix and attacker as in \(J_1\), retain
  \((M^{(3)},h_3)\), discard the original oracle and prefix views, and
  independently prepare
  \[
    V'_{\cA,3}=(W'_{\cA,3},S'_{\cA,3})
      \leftarrow\omega^A_{M^{(3)},h_3},
    \qquad
    V'_{\cB,3}=(W'_{\cB,3},S'_{\cB,3})
      \leftarrow\omega^B_{M^{(3)},h_3},
  \]
  set \(T'_{\cA,3}=T_{\cA}(W'_{\cA,3})\) and
  \(T'_{\cB,3}=T_{\cB}(W'_{\cB,3})\), and sample
  \(\widehat H_3\leftarrow\Unif(\cH(h_3,T'_{\cB,3}))\).
  \item Set \(\widetilde H_3=\widehat H_3[T'_{\cA,3}]\).
  \item \mrchanged{Run
  \((\Psi_3,K_{\cB,3})\leftarrow
    \cB_{\mathsf q}^{\widehat H_3}
    (1^\lambda,M^{(3)},S'_{\cB,3})\).}
  \item Run
  \(K^*_3\leftarrow
    \cA_{\mathsf{fin}}^{\widetilde H_3}
    (1^\lambda,M^{(3)},S'_{\cA,3},\Psi_3)\).
\end{hybrid}

\begin{lemma}[\(J_2\) to \(J_3\)]
\label{lem:multi-j2-j3}
\[
  \left|
    \Pr[J_2\text{ succeeds}]-\Pr[J_3\text{ succeeds}]
  \right|
  \le2q_{\cB,\mathsf q}\sqrt\eta.
\]
\end{lemma}

\begin{proof}
Condition on a reachable pair \((m,h)\) and on a value
\(W'_{\cB}=w'_{\cB}\) in the support of the classical-history marginal of
\(\omega^B_{m,h}\). Prepare \(S'_{\cB}\) in the corresponding conditional
state and set \(T'_{\cB}=T_{\cB}(w'_{\cB})\).
For the random reprogramming game of Katz and Sela
\cite[Lemma~3]{KatzSela24}, sample
\(\widehat H\leftarrow\Unif(\cH(h,T'_{\cB}))\) and set
\(F_0=\widehat H\). The randomized classical procedure \(C_{m,h}\) samples
the history \(W'_{\cA}\) from the classical marginal of
\(\omega^A_{m,h}\), sets \(T'_{\cA}=T_{\cA}(W'_{\cA})\), and returns
\[
  R
  =
  T'_{\cA}
  \big|_{\Dom(T'_{\cA})\setminus\Dom(h)}.
\]
Given challenge access to \(F_b\), run
\[
  (\Psi,K_{\cB})
  \leftarrow
  \cB_{\mathsf q}^{F_b}(1^\lambda,m,S'_{\cB}).
\]
After challenge access is removed, the coins of \(C_{m,h}\) reveal
\(W'_{\cA}\). Prepare \(S'_{\cA}\) in the corresponding conditional
state specified by \(\omega^A_{m,h}\), and run
\[
  K^*
  \leftarrow
  \cA_{\mathsf{fin}}^{\widehat H[T'_{\cA}]}
  (1^\lambda,m,S'_{\cA},\Psi).
\]
The distinguisher outputs \(1\) iff \(K^*=K_{\cB}\neq\bot\).

Every sampled Alice transcript is consistent with \(h\), so
\(F_1=\widehat H[R]=\widehat H[T'_{\cA}]\). Thus \(b=1\) reproduces \(J_2\)
and \(b=0\) reproduces \(J_3\), conditioned on
\((M,h,W'_{\cB})=(m,h,w'_{\cB})\). Only Bob queries the challenge oracle.
Moreover, by
\Cref{lem:multi-independence-attacker}, for every \(x\),
\[
  \Pr[x\in\Dom(R)]
  \le
  \begin{cases}
    0,&x\in\Dom(h),\\
    \Pr[x\in\Dom(T'_{\cA})]\le\eta,&x\notin\Dom(h).
  \end{cases}
\]
Thus \Cref{lem:reprogramming} bounds the absolute difference between the two
conditional success probabilities by \(2q_{\cB,\mathsf q}\sqrt\eta\) for
every \((m,h,w'_{\cB})\) in the support of the shared conditioning data. If
\(\pi(m,h,w'_{\cB})\) denotes the common subprobability mass of that data on
nonabort executions of \(J_2\) and \(J_3\), define
\[
  s_i(m,h,w'_{\cB})
  :=
  \Pr[J_i\text{ succeeds}\mid
    M^{(i)}=m,\ h_i=h,\ W'_{\cB,i}=w'_{\cB}]
  \qquad (i\in\{2,3\}).
\]
Because success is defined to be false on abort,
\[
  \Pr[J_i\text{ succeeds}]
  =
  \sum_{(m,h,w'_{\cB})\in\Supp(\pi)}
    \pi(m,h,w'_{\cB})s_i(m,h,w'_{\cB})
  \qquad (i\in\{2,3\}).
\]
The triangle inequality gives
\[
  \begin{aligned}
  \left|\Pr[J_2\text{ succeeds}]-\Pr[J_3\text{ succeeds}]\right|
  &\le
  \sum_{(m,h,w'_{\cB})\in\Supp(\pi)}\pi(m,h,w'_{\cB})
  \left|s_2(m,h,w'_{\cB})-s_3(m,h,w'_{\cB})\right|\\
  &\le 2q_{\cB,\mathsf q}\sqrt\eta.
  \end{aligned}
\]
\end{proof}

\begin{hybrid}{Multi-round hybrid \(J_4\)}
  \item Run the real prefix and attacker as in \(J_1\), retain
  \((M^{(4)},h_4)\), discard the original oracle and prefix views, and
  independently prepare
  \[
    V^{\mathsf E}_{\cA,4}
      =(W^{\mathsf E}_{\cA,4},S^{\mathsf E}_{\cA,4})
      \leftarrow\omega^A_{M^{(4)},h_4},
    \qquad
    V'_{\cB,4}=(W'_{\cB,4},S'_{\cB,4})
      \leftarrow\omega^B_{M^{(4)},h_4},
  \]
  set \(T^{\mathsf E}_{\cA,4}=T_{\cA}(W^{\mathsf E}_{\cA,4})\) and
  \(T'_{\cB,4}=T_{\cB}(W'_{\cB,4})\), and sample
  \(\widehat H_4\leftarrow\Unif(\cH(h_4,T'_{\cB,4}))\).
  \item \mrchanged{Independently prepare the Alice view that determines
  Bob's oracle,
  \[
    V^{\mathsf B}_{\cA,4}
    =
    (W^{\mathsf B}_{\cA,4},S^{\mathsf B}_{\cA,4})
    \leftarrow\omega^A_{M^{(4)},h_4},
  \]
  and set
  \(T^{\mathsf B}_{\cA,4}=T_{\cA}(W^{\mathsf B}_{\cA,4})\).
  Set
  \[
    H^{\mathsf B}_4=\widehat H_4[T^{\mathsf B}_{\cA,4}],
    \qquad
    H^{\mathsf E}_4=H^{\mathsf B}_4[T^{\mathsf E}_{\cA,4}].
  \]}
  \item Run
  \((\Psi_4,K_{\cB,4})\leftarrow
    \cB_{\mathsf q}^{H^{\mathsf B}_4}
    (1^\lambda,M^{(4)},S'_{\cB,4})\).
  \item Run
  \(K^*_4\leftarrow
    \cA_{\mathsf{fin}}^{H^{\mathsf E}_4}
    (1^\lambda,M^{(4)},S^{\mathsf E}_{\cA,4},\Psi_4)\).
\end{hybrid}

\begin{lemma}[\(J_3\) to \(J_4\)]
\label{lem:multi-j3-j4}
\[
  \left|
    \Pr[J_3\text{ succeeds}]-\Pr[J_4\text{ succeeds}]
  \right|
  \le
  2(q_{\cB,\mathsf q}+q_{\cA,\mathsf{fin}})\sqrt\eta.
\]
\end{lemma}

\begin{proof}
Condition on a reachable pair \((m,h)\) and on a value
\(W'_{\cB}=w'_{\cB}\) in the support of the classical-history marginal of
\(\omega^B_{m,h}\). Prepare \(S'_{\cB}\) in the corresponding conditional
state, and write \(T'_{\cB}=T_{\cB}(w'_{\cB})\). Independently prepare Eve's Alice view
\(V^{\mathsf E}_{\cA}=(W^{\mathsf E}_{\cA},S^{\mathsf E}_{\cA})
\leftarrow\omega^A_{m,h}\), and set
\(T^{\mathsf E}_{\cA}=T_{\cA}(W^{\mathsf E}_{\cA})\).  Sample
\(\widehat H\leftarrow\Unif(\cH(h,T'_{\cB}))\), output
\(F_0=\widehat H\), and let the randomized classical procedure
\(C_{m,h}\) independently sample \(W^{\mathsf B}_{\cA}\) from the classical
marginal of \(\omega^A_{m,h}\), set
\(T^{\mathsf B}_{\cA}=T_{\cA}(W^{\mathsf B}_{\cA})\), and return
\[
  R
  =
  T^{\mathsf B}_{\cA}
  \big|_{\Dom(T^{\mathsf B}_{\cA})\setminus\Dom(h)}.
\]
Given challenge access to \(F_b\), run Bob's final algorithm with \(F_b\)
and Alice's final algorithm with \(F_b[T^{\mathsf E}_{\cA}]\).  A query to
\(F_b[T^{\mathsf E}_{\cA}]\) is implemented with one query to \(F_b\) by
reversibly testing membership in the known transcript and supplying the
recorded answer on member inputs. The total number of challenge queries is
therefore at most
\(q_{\cB,\mathsf q}+q_{\cA,\mathsf{fin}}\).

After challenge access is removed, the revealed coins of \(C_{m,h}\)
determine \(W^{\mathsf B}_{\cA}\). The reduction can then prepare the
corresponding register \(S^{\mathsf B}_{\cA}\). This register is unused in
the \(b=0\) case and may be discarded; in the \(b=1\) case it completes the
auxiliary classical-quantum view sampled in \(J_4\). Every sampled Alice
transcript is consistent with \(h\). Therefore, for \(b=0\) the experiment
is \(J_3\) after discarding that unused view, while for \(b=1\),
\(F_1=\widehat H[T^{\mathsf B}_{\cA}]\) and the experiment is \(J_4\).
As in the preceding proof, each point belongs to the reprogramming set with
probability at most \(\eta\): points in \(\Dom(h)\) are excluded, and
Item~(iii) of \Cref{lem:multi-independence-attacker} applies everywhere else.
The random reprogramming lemma of Katz and Sela
\cite[Lemma~3]{KatzSela24} therefore bounds the absolute difference
between the conditional success probabilities of \(J_3\) and \(J_4\), given
\((M,h,W'_{\cB})=(m,h,w'_{\cB})\), by
\(2(q_{\cB,\mathsf q}+q_{\cA,\mathsf{fin}})\sqrt\eta\). Let
\(\pi(m,h,w'_{\cB})\) denote the common subprobability of these conditioning
data in the two hybrids, and define
\[
  s_i(m,h,w'_{\cB})
  :=
  \Pr[J_i\text{ succeeds}\mid
    M^{(i)}=m,\ h_i=h,\ W'_{\cB,i}=w'_{\cB}]
  \qquad (i\in\{3,4\}).
\]
The law of total probability and the triangle inequality yield
\[
  \begin{aligned}
  \left|\Pr[J_3\text{ succeeds}]-\Pr[J_4\text{ succeeds}]\right|
  &\le
  \sum_{m,h,w'_{\cB}}\pi(m,h,w'_{\cB})
  \left|s_3(m,h,w'_{\cB})-s_4(m,h,w'_{\cB})\right|\\
  &\le
  2(q_{\cB,\mathsf q}+q_{\cA,\mathsf{fin}})\sqrt\eta.
  \end{aligned}
\]
Again, the omitted abort branch has success probability zero in both
hybrids.
\end{proof}

\begin{hybrid}{Multi-round hybrid \(J_5\) (the real interception attack)}
  \item \mrchanged{Sample \(H_5\leftarrow\cH\), and run the honest classical prefix
  together with the online attacker.  Obtain
  \[
    (M^{(5)},h_5,V_{\cA,5},V_{\cB,5}),
  \]
  and abort if the attacker aborts.}
  \item Independently prepare Eve's Alice prefix view
  \[
    V'_{\cA,5}
    =(W'_{\cA,5},S'_{\cA,5})
    \leftarrow\omega^A_{M^{(5)},h_5},
  \]
  and set \(T'_{\cA,5}=T_{\cA}(W'_{\cA,5})\).
  \item \mrchanged{Run
  \((\Psi_5,K_{\cB,5})\leftarrow
    \cB_{\mathsf q}^{H_5}
    (1^\lambda,M^{(5)},S_{\cB,5})\),
  and transfer \(\Psi_5\) to Eve.}
  \item Run
  \[
    K^*_5
    \leftarrow
    \cA_{\mathsf{fin}}^{H_5[T'_{\cA,5}]}
    (1^\lambda,M^{(5)},S'_{\cA,5},\Psi_5).
  \]
\end{hybrid}

Fix a reachable pair \((m,h)\). The tuple
\[
  \bigl(
    H^{\mathsf B}_4,
    V^{\mathsf B}_{\cA,4},
    V'_{\cB,4}
  \bigr)
\]
in \(J_4\) has exactly the distribution of the conditioned prefix-resampling
experiment from \Cref{lem:prefix-resampling}. Eve's second Alice view
\(V^{\mathsf E}_{\cA,4}\) is independent in both hybrids. Hence, by
\Cref{lem:prefix-resampling} and trace-distance data processing
\cite[Chapter~3]{Watrous18}, for every reachable nonabort pair \((m,h)\),
\[
  \left|
    \Pr[J_4\text{ succeeds}\mid M^{(4)}=m,\ h_4=h]
    -
    \Pr[J_5\text{ succeeds}\mid M^{(5)}=m,\ h_5=h]
  \right|
  \le 2\eta.
\]
Let \(\pi(m,h)\) be the common probability that the prefix and attacker in
the two hybrids output \((m,h)\) without aborting, and define
\[
  s_i(m,h)
  :=
  \Pr[J_i\text{ succeeds}\mid M^{(i)}=m,\ h_i=h]
  \qquad (i\in\{4,5\}).
\]
The abort branch contributes zero, so the law of total probability and the
triangle inequality give
\[
  \left|
    \Pr[J_4\text{ succeeds}]-\Pr[J_5\text{ succeeds}]
  \right|
  \le
  \sum_{m,h}\pi(m,h)\left|s_4(m,h)-s_5(m,h)\right|
  \le 2\eta.
\]
Thus
\begin{equation}
\label{eq:multi-j4-j5}
  \left|
    \Pr[J_4\text{ succeeds}]-\Pr[J_5\text{ succeeds}]
  \right|
  \le2\eta.
\end{equation}

Hybrid \(J_5\) is precisely
\(\IntKR^{\mathsf{mr}}_{\Pi^{\mathsf{mr}},\cE_\eta}\). Combining
\eqref{eq:multi-j0-j1}, \eqref{eq:multi-j1-j2},
\Cref{lem:multi-j2-j3,lem:multi-j3-j4}, and
\eqref{eq:multi-j4-j5} gives
\begin{align*}
  \Adv^{\IntKR,\mathsf{mr}}_{
    \Pi^{\mathsf{mr}},\cE_\eta}(\lambda)
  &\ge
  \alpha
  -\eta-2\eta
  -2q_{\cB,\mathsf q}\sqrt\eta
  \\
  &\quad
  -2(q_{\cB,\mathsf q}+q_{\cA,\mathsf{fin}})\sqrt\eta
  -2\eta
  \\
  &=
  \alpha
  -5\eta
  -2(2q_{\cB,\mathsf q}+q_{\cA,\mathsf{fin}})\sqrt\eta.
\end{align*}
The query bound follows directly from
\Cref{lem:multi-independence-attacker,def:multi-eve}.
\end{proof}

\begin{corollary}[Polynomial-query multi-round impossibility]
\label{cor:multi-impossibility}
Let \(\Pi^{\mathsf{mr}}\) satisfy \Cref{def:multi-protocol}, let
\(\alpha=\alpha_{\Pi^{\mathsf{mr}}}(\lambda)>0\), and define
\[
  Q_{\mathsf{fin}}
  :=
  \max\{1,2q_{\cB,\mathsf q}+q_{\cA,\mathsf{fin}}\},
  \qquad
  \eta
  :=
  \left(\frac{\alpha}{16Q_{\mathsf{fin}}}\right)^2.
\]
Then the attacker of \Cref{thm:multi-main} succeeds with probability greater
than \(\alpha/2\) and makes
\[
  O\!\left(
    \frac{
      \max\{1,q_{\cA,\mathsf{pre}}\}
      \max\{1,q_{\cB,\mathsf{pre}}\}
      Q_{\mathsf{fin}}^6}
      {\alpha^6}
    +q_{\cA,\mathsf{fin}}
  \right)
\]
queries. It follows that any family of such protocols with
inverse-polynomial valid agreement and honest query complexity bounded by
\(\poly(\lambda)\) is insecure against a computationally unbounded quantum
attacker making at most \(\poly(\lambda)\) oracle queries.
The query bound depends on the number of prefix rounds only through the total
honest prefix-query bounds.
\end{corollary}

\begin{proof}
The choice of \(\eta\) satisfies \(\eta\le1/256<1/100\), and the
reprogramming loss is
\(2Q_{\mathsf{fin}}\sqrt\eta=\alpha/8\).  Since
\(Q_{\mathsf{fin}}\ge1\) and \(0<\alpha\le1\),
\(5\eta\le5\alpha/256\).  Thus the success lower bound is strictly greater
than \(\alpha/2\).  Substituting
\(1/\eta^3=(16Q_{\mathsf{fin}}/\alpha)^6\) gives the query bound.
\end{proof}

\section{Applications}\label{sec:corollaries}

In this section, we state the two main applications of our results.


\begin{definition}[Fully classical Alice and classical reply]
A protocol from \Cref{def:protocol} has fully classical Alice and a classical
reply if:
\begin{enumerate}[leftmargin=2em]
  \item \(\cA_1\) is a probabilistic classical oracle algorithm and its
  residual state is classical;
  \item Bob's message \(\Psi\) is classical; and
  \item \(\cA_2\) is a probabilistic classical algorithm with classical
  oracle access.
\end{enumerate}
Bob may still perform quantum computation and make quantum random-oracle
queries.
\end{definition}

\begin{corollary}[Classical eavesdropping attack]\label{cor:classical}
For a protocol with fully classical Alice and a classical reply, the
attacker in \Cref{thm:main} is a computationally unbounded classical
attacker with classical random-oracle access.  Therefore the conclusions of
\Cref{cor:parameters,cor:quantum-impossibility} hold against classical
eavesdroppers. Because the reply is classical, Eve can copy it and forward
it to Alice; the attack is therefore passive.
\end{corollary}

\begin{proof}
Under the stated assumptions, \(\omega_{m_1,h}\) is an ordinary classical
distribution over Alice's transcript and residual classical state.  The
attacker is classical, exact conditional sampling is classical, the reply can
be stored classically, and the simulation of \(\cA_2^{H[T'_{\cA}]}\) uses
only classical oracle queries.  No quantum operation remains in Eve's
algorithm.
\end{proof}

\begin{theorem}[Attack on QPKE with classical-query key generation]
\label{thm:qpke}
Let \(\Sigma\) satisfy \Cref{def:qpke}.  For every \(0<\delta<1\), there is
a computationally unbounded quantum IND-CPA attacker \(\cD_\delta\) making
at most
\[
  \ceil{q_{\Gen}/\delta^2}+q_{\Dec}
\]
random-oracle queries and satisfying
\[
  \Adv^{\IND\text{-}\mathrm{CPA}}_{\Sigma,\cD_\delta}(\lambda)
  \ge
  \alpha_\Sigma-\frac12-\delta
  -2(2q_{\Enc}+q_{\Dec})\sqrt\delta.
\]
In particular, if \(\alpha_\Sigma(\lambda)>0\), let
\[
  Q_\Sigma:=\max\{1,2q_{\Enc}+q_{\Dec}\},
  \qquad
  \delta:=\left(\frac{\alpha_\Sigma}{8Q_\Sigma}\right)^2.
\]
Then
\[
  \Adv^{\IND\text{-}\mathrm{CPA}}_{\Sigma,\cD_\delta}(\lambda)
  \ge \frac{47}{64}\alpha_\Sigma-\frac12,
\]
and \(\cD_\delta\) makes at most
\[
  \ceil{4096q_{\Gen}Q_\Sigma^4/\alpha_\Sigma^4}+q_{\Dec}
\]
random-oracle queries. Consequently, a scheme with negligible correctness
error and honest query complexity bounded by \(\poly(\lambda)\) is not
query-bounded IND-CPA secure: for every polynomial \(p\), an attacker making
at most \(\poly(\lambda)\) queries achieves advantage at least
\(1/2-1/p(\lambda)-\negl(\lambda)\), for every message-length function \(\mu\)
bounded above by a polynomial in \(\lambda\).
Encryption,
decryption, and the ciphertext may all be quantum. For
\(\mu(\lambda)=1\), this conclusion applies in particular to the
imperfectly correct PKE obtained from two-round OSP by Bartusek and
Khurana~\cite[Theorem~7.9]{BartusekKhurana24}, whenever the OSP sender's
random-oracle queries are classical.
\end{theorem}

\begin{proof}
Define a key-agreement protocol \(\Pi_\Sigma\) from \(\Sigma\) as follows.
Alice's first stage runs
\[
  (\pk,\sk)\leftarrow\Gen^H(1^\lambda)
\]
and outputs \((T_{\Gen},M_1,S_{\cA})=(T_{\Gen},\pk,\sk)\).  On input
\(\pk\), Bob fixes the two distinct messages
\[
  m_0:=0^{\mu(\lambda)},
  \qquad
  m_1:=0^{\mu(\lambda)-1}1,
\]
samples \(b\leftarrow\bits\), runs
\(C\leftarrow\Enc^H(1^\lambda,\pk,m_b)\), sends \(C\), and sets
\(K_{\cB}=b\). Alice's final stage computes
\(\widehat m\leftarrow\Dec^H(1^\lambda,\sk,C)\) and outputs
\[
  K_{\cA}:=
  \begin{cases}
    0 & \text{if }\widehat m=m_0,\\
    1 & \text{if }\widehat m=m_1,\\
    \bot & \text{otherwise.}
  \end{cases}
\]
For \(b\in\bits\), let
\begin{align*}
  p_b:=\Pr\bigl[&\widehat m=m_b:\ \\
    &H\leftarrow\cH_\lambda,\
    (\pk,\sk)\leftarrow\Gen^H(1^\lambda),\\
    &C_b\leftarrow\Enc^H(1^\lambda,\pk,m_b),\
    \widehat m\leftarrow\Dec^H(1^\lambda,\sk,C_b)\bigr].
\end{align*}
Therefore
\[
  q_{\cA,1}=q_{\Gen},\qquad
  q_{\cB}=q_{\Enc},\qquad
  q_{\cA,2}=q_{\Dec},
  \qquad
  \alpha_{\Pi_\Sigma}=\frac{p_0+p_1}{2}\ge\alpha_\Sigma,
\]
where the inequality follows from the definition of the minimum
correct-decryption probability in \Cref{def:qpke}.

This is the standard reduction from QPKE to key agreement used in prior
QROM impossibility work~\cite{BouazizEtAl24,LiEtAl24}. Apply
\Cref{thm:main} to obtain the
intercepting attacker
\(\cE_\delta=(\cE_{1},\cE_{2})\).  The IND-CPA
attacker \(\cD_\delta\) has its first stage run
\[
  Z_{\cD}\leftarrow\cE_{1}^H(1^\lambda,\pk)
\]
and returns \((Z_{\cD},m_0,m_1)\). Its second stage, upon
receiving the challenge ciphertext register \(C\), runs
\(K^*\leftarrow\cE_{2}^H(1^\lambda,Z_{\cD},C)\) and outputs \(K^*\)
when \(K^*\in\bits\), or an independent uniform bit otherwise.  Thus its
success probability satisfies
\[
  \Pr[\IND_{\Sigma,\cD_\delta}(1^\lambda)=1]
  \ge
  \Pr[\IntKR_{\Pi_\Sigma,\cE_\delta}(1^\lambda)=1].
\]
The claimed advantage and query bound follow from \Cref{thm:main}, since
\begin{align*}
  \Adv^{\IND\text{-}\mathrm{CPA}}_{\Sigma,\cD_\delta}(\lambda)
  &\ge
  \Pr[\IND_{\Sigma,\cD_\delta}(1^\lambda)=1]-\frac12\\
  &\ge
  \Pr[\IntKR_{\Pi_\Sigma,\cE_\delta}(1^\lambda)=1]-\frac12.
\end{align*}

For the stated parameter choice,
\[
  2Q_\Sigma\sqrt\delta=\frac{\alpha_\Sigma}{4},
  \qquad
  \delta\le\frac{\alpha_\Sigma}{64},
\]
which gives the factor \(47/64\); the query bound follows by substituting
\(1/\delta^2=4096Q_\Sigma^4/\alpha_\Sigma^4\).  A quantum challenge
ciphertext is transferred to the attacker's second stage, so the reduction
does not copy it.  If \(\Gen\) and \(\Dec\) are classical algorithms and the
ciphertext is classical, \Cref{cor:classical} moreover yields a classical
attack.
\end{proof}

\bibliographystyle{alpha}
\bibliography{quantum}


\end{document}
